\chapter{场线、等势面与多极场}\label{ch:es1}

静电场既可由电荷直接叠加，也可由电势的梯度求得。场线给出电场方向，等势面给出电势相同的位置。讨论通量、零场点和远处形状时，需要把图中的几何关系写进场方程。

点电荷与无限线电荷的场线具有不同几何；有限电荷对的展开则给出偶极和四极近似。使用远区式时，观察点须远离整个源区，角度按图中所标的轴量取。全书采用SI单位，真空中的库仑常数记为$\ke=1/(4\pi\epsilon_0)$。


\section{源、场与电势}\label{sec:es1-foundation}
有限电荷分布可直接求场，也可先求标量电势再求导。两条路线均须区分源坐标与观察坐标，明确零势，并判断对称性实际确定了哪些场分量。

\subsection{电荷分布与叠加}
点电荷$Q$在$\vct r_0$产生的场与势为
\[\vct E(\vct r)=\ke Q\frac{\vct r-\vct r_0}{|\vct r-\vct r_0|^3},\qquad
\varphi(\vct r)=\frac{\ke Q}{|\vct r-\vct r_0|}.\]
电势的零点取在无穷远，要求源局限于有限区域。连续分布则将求和变成积分，按线、面、体分布分别使用$\lambda\dd l$、$\sigma\dd S$、$\rho\dd V$，不可把不同量纲的密度混用。体分布的势为
\[\varphi(\vct r)=\ke\int\frac{\rho(\vct r')}{|\vct r-\vct r'|}\dd V'.\]
积分变量$\vct r'$属于源，观察点$\vct r$在求导时变化。场的叠加是矢量相加，势的叠加是标量相加；分量运算中的方向符号应由坐标与电荷符号共同确定。

电势$\varphi$、密度$\rho$是标量场；电场$\vct E$、磁感应强度$\vct B$、电流密度$\vct j$是矢量场。写成$\vct A=A_x\vct e_x+A_y\vct e_y+A_z\vct e_z$时，三个分量仍是随位置变化的标量函数。在曲线坐标中，单位基矢本身也可能随位置变化。

\subsection{积分关系与局部关系}
静电场满足
\[\iint_{\partial V}\vct E\cdot\vct n\dd S=\frac{Q_V}{\epsilon_0},\qquad
\oint_C\vct E\cdot\dd\vct l=0.\]
第一个积分使用闭曲面的外法向，计入内部全部电荷；第二个积分的路径须处于静电场中。对有限电荷分布，任何不穿过源的封闭面都可使用Gauss定理，但仅有通量还不足以确定曲面上各点的场，除非另有对称性。

在光滑区域内，积分式分别对应
\[\nabla\cdot\vct E=\frac\rho{\epsilon_0},\qquad
\nabla\times\vct E=0.\]
由静电环路性质定义单值电势，得$\vct E=-\nabla\varphi$。代入散度关系，
\[\boxed{\nabla^2\varphi=-\frac\rho{\epsilon_0}}.\]
无源处才可化为Laplace方程。点电荷处的场并不光滑，不能把原点以外的$\nabla\cdot\vct E=0$用到包围该点的整个体积。

\subsection{场线、等势面与实际轨迹}
场线在每个普通点的切向与$\vct E$平行，其参数方程可写作$\dd\vct r/\dd u=f(u)\vct E$。在二维截面内且$E_x\ne0$时，$\dd y/\dd x=E_y/E_x$；分母为零的竖直切线应改换参数。零场点处没有唯一的场线切向，需另分析邻近场。

等势面满足$\varphi=C$。沿面移动有$\dd\varphi=\nabla\varphi\cdot\dd\vct r=0$，故场在非零处垂直等势面。这里的正交是实际空间中的正交；若横纵轴使用不同绘图比例，纸面上看出的夹角未必保持直角。

\begin{example}[两种方向场的几何]\label{ex:es1-e1}
平面内$\varphi(x,y)=Axy$，$A>0$。求电场、等势线与场线，并说明带电粒子从静止释放后是否必沿场线运动。 见图~\ref{fig:es1-e1-q}。
\end{example}

\begin{figure}[H]\centering
\begin{tikzpicture}[>=Stealth,line cap=round,line join=round]
\begin{scope}[x=1cm,y=1cm]\draw[axis](0,0)--(2.18,0)node[right]{$x$};\draw[axis](0,0)--(0,2.1)node[above]{$y$};\draw[curve,domain=0.4:2,samples=101]plot(\x,{.7/\x});\node[right]at(2,.8){$xy=0.7$};\end{scope}

\end{tikzpicture}
\caption{第一象限内的一条等势线，横纵坐标采用相同长度单位。}\label{fig:es1-e1-q}
\end{figure}
\begin{solution}
偏导时另一个坐标保持不变，故$E_x=-Ay$、$E_y=-Ax$。等势线为$xy=C/A$，包括零势时的两条坐标轴。场线方程为
\[\begin{aligned}
\frac{\dd y}{\dd x}&=\frac xy,\qquad y\dd y=x\dd x,\\
\int y\dd y&=\int x\dd x,\qquad \frac{y^2}2=\frac{x^2}2+C_0,
\quad \boxed{y^2-x^2=C_1}.
\end{aligned}\]
在$y=0$处不能直接使用含$1/y$的斜率式，但隐式曲线仍给出正确的切向。$y=\pm x$是经过零场点的特殊分界方向，原点本身$\vct E=0$。

图~\ref{fig:es1-e1-s}的两族双曲线满足$\nabla(xy)\cdot\nabla(y^2-x^2)=0$，所以在普通交点相互正交。每条双曲线都有两支，图中应按同一常数保留它们。

粒子的加速度与电场平行，速度却由此前各时刻的加速度累积。曲线运动需要法向加速度；若速度严格沿一条弯曲场线，而电力始终切向，就不能提供该曲率所需的法向分量。因此一般轨迹不会与场线重合。从静止释放只能确定最初加速方向，不能据此把整段轨迹当作场线。
\begin{figure}[H]\centering
\begin{tikzpicture}[>=Stealth,line cap=round,line join=round]

\draw[axis](-2,0)--(2.2,0)node[right]{$x$};\draw[axis](0,-2)--(0,2.2)node[above]{$y$};
\foreach\c in{.4,.8}{\draw[curve,domain=.4:1.8,samples=81]plot(\x,{\c/\x});\draw[curve,domain=-1.8:-.4,samples=81]plot(\x,{\c/\x});}
\draw[dashed,domain=-1.6:1.6,samples=81]plot(\x,{sqrt((\x)^2+.45)});\draw[dashed,domain=-1.6:1.6,samples=81]plot(\x,{-sqrt((\x)^2+.45)});
\node[right,align=left]at(2.7,.6){实线：$xy=C$\\[5pt]虚线：$y^2-x^2=C_1$};


\end{tikzpicture}
\caption{等势线与场线使用同一坐标比例；所画实线取正的等势常数。}\label{fig:es1-e1-s}
\end{figure}


若还要写出实际运动，可令$u=x+y$、$v=x-y$。对正电荷记$\omega=\sqrt{qA/m}$，运动方程变为$\ddot u+\omega^2u=0$、$\ddot v-\omega^2v=0$。特征根分别为$\pm\mathrm i\omega$与$\pm\omega$，通解为
\[u=A_1\cos\omega t+B_1\sin\omega t,\qquad
v=A_2e^{\omega t}+B_2e^{-\omega t}.\]
从$(x_0,y_0)$静止释放时，$u(0)=x_0+y_0$、$v(0)=x_0-y_0$、$\dot u(0)=\dot v(0)=0$，因此$A_1=x_0+y_0$、$B_1=0$、$A_2=B_2=(x_0-y_0)/2$。由$x=(u+v)/2$、$y=(u-v)/2$恢复轨迹。一般初值同时激发振荡模和指数模；只有特殊直线方向上的初值才可能使轨迹沿直场线。场线方程只规定局部切向，牛顿方程还需要初速度。
\end{solution}
\subsection{电场管与有向通量}
取由一束场线围成的细管。管侧面与电场相切，侧面通量为零。两个截面的通量若都按沿管方向计为$\Phi_1,\Phi_2$，闭曲面的有向通量是$\Phi_2-\Phi_1$，而不是二者之和。截面法向与管方向相反时，其符号随之改变。

\begin{example}[截面通量何时守恒]\label{ex:es1-e2}
一条电场管在两个截面之间包围电荷$q$。两截面的面积分别为$S_1,S_2$，电场近似均匀且垂直截面。沿场线方向记录场强的有符号分量$E_1,E_2$。求它们的关系，讨论无源与有源两种情况。 见图~\ref{fig:es1-e2-q}。
\end{example}

\begin{figure}[H]\centering
\begin{tikzpicture}[>=Stealth,line cap=round,line join=round]

\draw[curve](0,.7)..controls(2,1.1)and(3,1.3)..(4.5,1.6);\draw[curve](0,-.7)..controls(2,-1.1)and(3,-1.3)..(4.5,-1.6);
\draw(0,0)ellipse(.18 and .7);\draw(4.5,0)ellipse(.25 and 1.6);\node[below]at(0,-.8){$S_1$};\node[below]at(4.5,-1.7){$S_2$};
\draw[axis](1.5,0)--(2.6,0);\fill(3,.25)circle(1.6pt)node[above]{$q$};

\end{tikzpicture}
\caption{侧面沿场线构成，截面通量均按从左向右记录。}\label{fig:es1-e2-q}
\end{figure}
\begin{solution}
图~\ref{fig:es1-e2-q}的侧壁上$\vct E\cdot\vct n=0$。在闭合管体的左端，外法向向左；右端外法向向右，故Gauss定理给出
\[\boxed{E_2S_2-E_1S_1=\frac q{\epsilon_0}}.\]
图~\ref{fig:es1-e2-s}把两端的外法向另行标出，避免把同向记录的通量当成两个都向外的通量。无源时$q=0$，有$E_1S_1=E_2S_2$，管束变宽处场强相应减小；有净电荷时，两截面通量一般不等。

若截面上场不均匀，记沿管方向的通量为$\Phi_i=\int_{S_i}\vct E\cdot\vct n_i^{\rightarrow}\dd S$，则应使用$\Phi_2-\Phi_1=q/\epsilon_0$。只有场近似均匀且垂直截面时，才能写成$E_iS_i$。图上的线密度要与场强成比例，还须预先规定各根细管携带相同通量；任意选画的几条线只能表示方向。

当管体穿过点电荷时，源点会使场线的起止结构改变，不能沿着同一无源管继续使用通量常数。后面用固体角求场线时，参考曲面跨越电荷也会产生相应的常数跳变。
\begin{figure}[H]\centering
\begin{tikzpicture}[>=Stealth,line cap=round,line join=round]

\draw[curve](0,-.7)rectangle(4.5,.7);\draw[axis](0,0)--(-.9,0)node[left]{$\vct n_1$};\draw[axis](4.5,0)--(5.4,0)node[right]{$\vct n_2$};
\node at(1,.25){$-E_1S_1$};\node at(3.4,.25){$+E_2S_2$};\node at(2.2,-.3){$q/\epsilon_0$};

\end{tikzpicture}
\caption{闭合管体的外法向在两端反向，通量差由管内净电荷决定。}\label{fig:es1-e2-s}
\end{figure}


若侧面未选成场线面，还须计入向外侧通量$\Phi_{\rm side}$，关系改为$\Phi_2-\Phi_1+\Phi_{\rm side}=q/\epsilon_0$。截面变宽使平均场减小的推断，因而需要无源、无侧向漏通量这两个条件。
\end{solution}

\begin{example}[两个异号点电荷的场线]\label{ex:es1-e3}
$+Q$位于$x=-l$，$-Q$位于$x=l$。求轴对称场的场线方程，以$xy$平面表示子午截面，$y\ge0$。用分量微分与通量两种途径说明所得曲线，并处理参考面改变时的通量常数。 见图~\ref{fig:es1-e3-q}。
\end{example}

\begin{figure}[H]\centering
\begin{tikzpicture}[>=Stealth,line cap=round,line join=round]
\draw[axis](-2.7,0)--(2.9,0)node[right]{$x$};\fill(-1.6,0)circle(1.6pt)node[below]{$+Q$};\fill(1.6,0)circle(1.6pt)node[below]{$-Q$};\fill(0.3,1.8)circle(1.5pt)node[above right]{$P$};\draw[guide](-1.6,0)--(0.3,1.8);\draw[guide](1.6,0)--(0.3,1.8);\node[above]at(-.7,.65){$r_+$};\node[above]at(1.05,1.05){$r_-$};
\end{tikzpicture}
\caption{两源在同一轴上，空间中的场线由子午曲线绕轴旋转得到。}\label{fig:es1-e3-q}
\end{figure}
\begin{solution}
令$r_\pm=\sqrt{(x\pm l)^2+y^2}$，其中加号下标对应左侧正电荷。分量叠加为
\[E_x=\ke Q\left(\frac{x+l}{r_+^3}-\frac{x-l}{r_-^3}\right),\qquad
E_y=\ke Qy\left(\frac1{r_+^3}-\frac1{r_-^3}\right).\]
因此场线斜率为$E_y/E_x$。与直接解这个分式方程相比，寻找一个沿场线不变的通量函数更方便。定义
\[F(x,y)=\frac{x+l}{r_+}-\frac{x-l}{r_-}.\]
例如对左源的一项应用乘积法则与链式法则，
\[\partial_x\frac{x+l}{r_+}=\frac1{r_+}-\frac{(x+l)^2}{r_+^3}=\frac{y^2}{r_+^3},\qquad
\partial_y\frac{x+l}{r_+}=-\frac{(x+l)y}{r_+^3}.\]
右源只需将$x+l,r_+$换成$x-l,r_-$，并保留$F$中的减号，故
\[F_x=\frac{yE_y}{\ke Q},\qquad F_y=-\frac{yE_x}{\ke Q}.
\]
于是沿$\dd y/\dd x=E_y/E_x$的曲线，$\dd F/\dd x=0$。所求隐式方程为
\[\boxed{\frac{x+l}{\sqrt{(x+l)^2+y^2}}
-\frac{x-l}{\sqrt{(x-l)^2+y^2}}=C}.\]

图~\ref{fig:es1-e3-s}将$P$与两电荷连接，记这两条射线同正$x$方向的夹角为$\theta_+,\theta_-$，则$F=\cos\theta_+-\cos\theta_-$。点电荷穿过固体角$\Omega$的通量为$Q\Omega/(4\pi\epsilon_0)$；绕轴的球冠固体角是$2\pi(1-\cos\theta)$。两源的球冠通量相加后，只差一个固定的参考常数，故$F$不变也等价于同一场管通量不变。

若换用另一个帽面或把参考截面移到负电荷的另一侧，两种闭合方式之间可能多包围一个$-Q$。此时通量常数相差$Q/\epsilon_0$，而表示实际曲线的角度关系不变。常数的数值取决于如何封口，不能在跨过源点后强行令两种参考下的常数相等。

轴线上的区间还须单列，$y=0$使上面的微分关系同时含零因子。在$x<-l$、$-l<x<l$和$x>l$三个区间，$E_x$的方向依次向左、向右和向左。中间一支由正电荷通向负电荷；左侧外支由正电荷伸向无穷远，右侧外支由无穷远通向负电荷。后两者是轴上极限分支，不能把它们也说成有限长度的两源连接线。轴外的普通场线由正源通向负源，图中的对称关系来自等量异号的源条件。
\begin{figure}[H]\centering
\begin{tikzpicture}[>=Stealth,line cap=round,line join=round]
\draw[axis](-2.7,0)--(2.9,0)node[right]{$x$};\fill(-1.6,0)circle(1.6pt)node[below]{$+Q$};\fill(1.6,0)circle(1.6pt)node[below]{$-Q$};\fill(0.3,1.8)circle(1.5pt)node[above right]{$P$};\draw[guide](-1.6,0)--(0.3,1.8);\draw[guide](1.6,0)--(0.3,1.8);
\draw(-.95,0)arc(0:43.45:.65);\node at(-.85,.4){$\theta_+$};
\draw(2.2,0)arc(0:125.84:.6);\node at(1.9,.65){$\theta_-$};
\draw[axis](.3,1.8)--(1.15,1.5575)node[above]{$\vct E$};

\end{tikzpicture}
\caption{角度都从正$x$方向量起，负电荷一侧的角可以大于$\pi/2$；不得将它换成补角却仍沿用原式。}\label{fig:es1-e3-s}
\end{figure}


在垂直平分线$x=0$上，$F=2l/\sqrt{l^2+y^2}$，随交点高度从$2$降到$0$。这给场线常数一个直接的几何编号，也说明不同高度的曲线不会在普通点相交。
\end{solution}

\begin{example}[异号平行线电荷的两族圆]\label{ex:es1-e4}
两条无限长直线电荷平行于$z$轴，在截面上位于$x=\pm l$，左线为$+\lambda$、右线为$-\lambda$。求截面内的场线和等势线，证明它们分别是两族圆，并讨论无穷远电势基准。 见图~\ref{fig:es1-e4-q}。
\end{example}

\begin{figure}[H]\centering
\begin{tikzpicture}[>=Stealth,line cap=round,line join=round]
\draw[axis](-2.7,0)--(2.9,0)node[right]{$x$};\fill(-1.5,0)circle(1.6pt)node[below]{$+\lambda$};\fill(1.5,0)circle(1.6pt)node[below]{$-\lambda$};\fill(0,1.8)circle(1.5pt)node[above right]{$P$};\draw[guide](-1.5,0)--(0,1.8);\draw[guide](1.5,0)--(0,1.8);
\end{tikzpicture}
\caption{沿$z$方向平移不改变场，问题化为二维截面。}\label{fig:es1-e4-q}
\end{figure}
\begin{solution}
单条无限线的场为$E=\lambda/(2\pi\epsilon_0r)$。沿径向积分，
\[\varphi(r_2)-\varphi(r_1)=-\int_{r_1}^{r_2}\frac{\lambda}{2\pi\epsilon_0r}\dd r
=-\frac{\lambda}{2\pi\epsilon_0}[\ln r]_{r_1}^{r_2}
=-\frac{\lambda}{2\pi\epsilon_0}\ln\frac{r_2}{r_1}.\]
单线不能把无穷远规定为有限零势，但两条异号等量线的对数参考长度相消，得到
\[\varphi=\frac\lambda{2\pi\epsilon_0}\ln\frac{r_-}{r_+},\qquad
r_\pm=\sqrt{(x\pm l)^2+y^2}.\]
该组合势在远处趋零。等势条件为$r_-/r_+=\kappa>0$，平方整理，
\[(1-\kappa^2)(x^2+y^2+l^2)-2l(1+\kappa^2)x=0.\]
$\kappa\ne1$时是圆，圆心在$x$轴上；$\kappa=1$退化为垂直平分线$x=0$。

场线可由场强分量求斜率，也可利用单位长度通量。两射线的夹角保持不变，圆周角定理表明场线是经过$(-l,0),(l,0)$的圆弧。在上半平面写成
\[\boxed{x^2+(y-b)^2=l^2+b^2},\qquad b\in\mathbb R.
\]
只取从正线指向负线的相应圆弧，不把源点处的奇异切线当作普通点。图~\ref{fig:es1-e4-s}给出一个过源圆与一个等势圆。

若过源圆的圆心为$(0,b)$，等势圆圆心为$(a,0)$、半径$R$，由等势圆方程有$a^2-R^2=l^2$。两圆圆心距平方$a^2+b^2$等于半径平方之和$R^2+(l^2+b^2)$，所以两圆正交。这是$\vct E=-\nabla\varphi$的具体几何表现。
\begin{figure}[H]\centering
\begin{tikzpicture}[>=Stealth,line cap=round,line join=round]

\draw[axis](-2.4,0)--(3.1,0)node[right]{$x$};\draw[axis](0,-1.7)--(0,2.6)node[above]{$y$};
\draw[curve](0,.65)circle(1.63554);\draw[dashed](2,0)circle(1.32288);
\fill(-1.5,0)circle(1.8pt)node[below]{$+\lambda$};\fill(1.5,0)circle(1.8pt)node[below]{$-\lambda$};
\node[above]at(-1.1,2){场线圆};\node[right]at(2.8,1){等势圆};

\end{tikzpicture}
\caption{实线圆通过两条源线的截点，虚线圆的两源距离比为常数；交点处相互正交。}\label{fig:es1-e4-s}
\end{figure}


两线的分量可写成$E_x=\lambda[(x+l)/r_+^2-(x-l)/r_-^2]/(2\pi\epsilon_0)$，$E_y=\lambda y(1/r_+^2-1/r_-^2)/(2\pi\epsilon_0)$。整理后$E_y/E_x=-2xy/(l^2-x^2+y^2)$。而过源圆的隐式导数为$y'=-x/(y-b)$，由圆方程消去$b$，恰得同一斜率。圆周角构造与分量法得到同一个斜率方程，因而确定相同的场线。
\end{solution}

\begin{example}[距离比与反演点]\label{ex:es1-e5}
圆$|\vct r|=R$上取任意点$A$，点$P$位于$x=d>R$处。求轴上点$P'$，使$|A-P'|/|A-P|$对整个圆周为常数，并由此构造一个球面的零势电荷对。 见图~\ref{fig:es1-e5-q}。
\end{example}

\begin{figure}[H]\centering
\begin{tikzpicture}[>=Stealth,line cap=round,line join=round]
\draw[curve](0,0)circle(1.4);\fill(0,0)circle(1.4pt)node[below left]{$O$};\fill(2.5,0)circle(1.6pt)node[below]{$P$};\fill(.784,0)circle(1.6pt)node[below]{$P\prime$};\fill(.7,1.21244)circle(1.6pt)node[above]{$A$};\draw[guide](.7,1.21244)--(2.5,0) (.7,1.21244)--(.784,0);
\end{tikzpicture}
\caption{反演点位于$OP$上，必须对全部球面点保持同一距离比例。}\label{fig:es1-e5-q}
\end{figure}
\begin{solution}
设$P'$在$x=d'$处，$A$的极角从$OP$量起，则
\[|A-P|^2=R^2+d^2-2Rd\cos\theta,
\quad |A-P'|^2=R^2+d'^2-2Rd'\cos\theta.\]
要求后式为前式的固定倍数，比较常数项与$\cos\theta$系数，得到非平凡解
\[\boxed{d'=R^2/d,\qquad\frac{|A-P'|}{|A-P|}=R/d}.\]
图~\ref{fig:es1-e5-s}所示三角形也可用相似关系证明这一比例，且对绕轴旋转后的整个球面仍成立。

若真电荷$Q$在$P$，在$P'$放置$Q'=-RQ/d$，球面上的势为
\[\ke\left(\frac Q{|A-P|}+\frac{Q'}{|A-P'|}\right)=0.\]
这已经给出第三讲接地导体球的像电荷组合。这里先建立的是球面上的等势关系；要把它用于物理边值问题，还必须确认像点在所求区域之外、真实源没有改变，并满足远场条件，再由唯一性判定。

反演距离与电量比例都是由同一个距离恒等式确定，不能把$d'=R^2/d$与任意选取的像电量拼在一起。若$P$在球内，反演点移到球外；相同恒等式依然成立，但物理区域与边界条件应随之改变。
\begin{figure}[H]\centering
\begin{tikzpicture}[>=Stealth,line cap=round,line join=round]
\draw[curve](0,0)circle(1.4);\draw[axis](-.4,0)--(3,0);
\fill(0,0)circle(1.4pt)node[below left]{$O$};
\fill(2.5,0)circle(1.6pt)node[below]{$P$};\fill(.784,0)circle(1.6pt)node[below]{$P'$};
\coordinate(A)at(.7,1.21244);\fill(A)circle(1.5pt)node[above]{$A$};
\draw(A)--(2.5,0)node[pos=.65,above right]{$AP$};
\draw(A)--(.784,0)node[pos=.55,left]{$AP'$};\draw[guide](A)--(0,0);
\node[right]at(3.15,.75){$P:Q$};\node[right]at(3.15,.15){$P':-RQ/d$};\node[right]at(3.15,-.65){$AP'/AP=R/d$};
\end{tikzpicture}
\caption{球面上两个距离之比固定，使两个异号电荷的势逐点抵消。}\label{fig:es1-e5-s}
\end{figure}


比较距离平方时，若记比例平方为$t$，先有$d'=td$和$R^2+d'^2=t(R^2+d^2)$，消去$t$得$(d-d')(R^2-dd')=0$。第一根$d'=d$只是同一点的平凡解，不能拿来放置另一个电荷；第二根才是反演关系。像点在球内而真实源在球外，正是外域唯一性所需的奇性分离。
\end{solution}

\subsection{由电势积分求场的步骤}
当观察点不在源分布上且积分收敛时，可将梯度移入势积分。令$\vct R=\vct r-\vct r'$，有$\nabla(1/R)=-\vct R/R^3$，所以
\[\vct E=-\nabla\varphi=\ke\int\rho(\vct r')\frac{\vct R}{R^3}\dd V'.\]
这既解释了势法与场法的等价，也说明求导不能同时移动源坐标。若电荷在一张曲面上，势可连续而法向导数有跃变；这时先在面两侧分别求导，再按边界取极限。直接把观察点放到带电面上积分，容易把自场奇性与外加场混在一起。

无限线电荷可先用长度为$2L$的有限线段说明参考势的问题。在垂直平分线上距线$s$处，
\[\begin{aligned}
\varphi_L(s)&=\ke\lambda\int_{-L}^L\frac{\dd z'}{\sqrt{s^2+z'^2}}\\
&=\ke\lambda[\arsinh(z'/s)]_{-L}^L=2\ke\lambda\arsinh(L/s)\\
&=2\ke\lambda\left[\ln(2L/s)+O(s^2/L^2)\right].
\end{aligned}
\]
这里$s>0$，最后一步使用
\[\begin{aligned}
\arsinh(L/s)&=\ln\frac{L+\sqrt{L^2+s^2}}s,\\
\sqrt{L^2+s^2}&=L\left[1+\frac{s^2}{2L^2}+O(s^4/L^4)\right],\\
\arsinh(L/s)&=\ln\frac{2L}s+
\ln\left[1+\frac{s^2}{4L^2}+O(s^4/L^4)\right]
=\ln\frac{2L}s+O(s^2/L^2).
\end{aligned}\]
固定参考距离$s_0>0$后，当$L\to\infty$，$\varphi_L(s)$含有与$s$无关的发散部分$2\ke\lambda\ln(2L/s_0)$，而势差$\varphi_L(s_2)-\varphi_L(s_1)$趋于$-2\ke\lambda\ln(s_2/s_1)$。对数势中的参考长度由这一极限确定。

\subsection{场线通量函数的微分结构}
轴对称电场取球坐标$(r,\theta)$，定义$\Psi=2\pi\int_0^\theta E_r(r,u)r^2\sin u\dd u$。若从极轴到该点的球冠不穿过源，散度关系给
\[\partial_\theta\Psi=2\pi r^2\sin\theta E_r,\qquad
\partial_r\Psi=-2\pi r\sin\theta E_\theta.
\]
后一式来自无源散度方程
\[\partial_r(r^2E_r)\sin u=-r\partial_u(\sin uE_\theta).\]
在场光滑且极轴处$E_\theta$有界的球冠内，固定$r$积分得到
\[\partial_r\Psi=-2\pi r\int_0^\theta\partial_u(\sin uE_\theta)\dd u
=-2\pi r[\sin uE_\theta]_0^\theta=-2\pi r\sin\theta E_\theta.\]
沿切向场线$\dd r/(r\dd\theta)=E_r/E_\theta$，两项相消，故$\dd\Psi=0$。参考积分区若跨越点源，积分常数须按源通量调整。

在柱坐标$(x,s)$的轴对称场中，也可定义$\Psi=2\pi\int_0^s E_x(x,u)u\dd u$，有$\Psi_s=2\pi sE_x$、$\Psi_x=-2\pi sE_s$。后面异号点电荷的不变量正是这一结构。二维无限线电荷的单位长度通量没有旋转面积因子$2\pi s$，所以场线方程会与点电荷不同。


\subsection{有限源展开的余项与观察区域}\label{sec:es-series-error}
设全部源点位于半径$a$的球内，观察距离$r>a$。对一个源点，令$\gamma$为$\vct r$与$\vct r'$的夹角，库仑核可写为
\[\frac1{|\vct r-\vct r'|}
=\frac1r\frac1{\sqrt{1-2t\cos\gamma+t^2}},\qquad t=\frac{r'}r<1.
\]
把根号倒数按$t$展开，前三级为
\[\frac1{|\vct r-\vct r'|}=\frac1r+
\frac{r'\cos\gamma}{r^2}+
\frac{r'^2(3\cos^2\gamma-1)}{2r^3}+\cdots.
\]
每一级都须先对全部源求和。总电量为零只消去第零级；再有总偶极矩为零，才从四极级开始。单个正负源的势可能很大，彼此抵消后留下的小量仍由这些矩控制，不能以某一电荷单独产生的势作为最后答案的相对误差分母。

完整核展开由生成函数
\[\frac1{\sqrt{1-2tc+t^2}}=\sum_{l=0}^{\infty}P_l(c)t^l,
\qquad |t|<1,\quad -1\le c\le1
\]
定义，其中$P_l$为勒让德多项式；前三个是$1,c,(3c^2-1)/2$。对离散源记$Q_{\rm abs}=\sum_i|q_i|$。利用实角度下$|P_l(\cos\gamma)|\le1$，截断到$l=L$后的势余项满足
\[\begin{aligned}
|\mathcal R_L|&\le\frac{\ke Q_{\rm abs}}r
\sum_{l=L+1}^{\infty}(a/r)^l\\
&=\frac{\ke Q_{\rm abs}}r
\frac{(a/r)^{L+1}}{1-a/r}.
\end{aligned}\]
连续源将$Q_{\rm abs}$换成$\int|\rho|\dd V$即可。这个界对整个外球面成立，包含首项角因子为零的方向。若首个非零多极势在某条射线上消失，绝对误差仍受上式控制，但相对误差未必小。求电场还要对余项作空间求导；在$r\ge r_0>a$的闭外域内逐项求导是合法的，每次求导多带一个距离倒数。

取理想偶极极限时，须先说明保持不变的物理量。固定$q$令间距$\ell\to0$，偶极矩$p=q\ell$也趋零；要得到非零的理想偶极，须令$q\ell=p$固定。在距源固定的区域，极限场可以有限，源区的自身能却不因此有限。远区展开要求观察点远离源；理想偶极极限则规定源电荷与间距怎样共同变化。



\section{场的微分与曲线坐标}\label{sec:es1-operators}
局部场方程的形式不依赖坐标，分量表达却包含几何尺度。先从有向线元、面积元和体积元建立算子，再把对称性用于具体场。

\subsection{梯度、散度与旋度的定义}
可微标量场的方向变化为$\dd f=\nabla f\cdot\dd\vct r$。梯度是这个线性变化的矢量系数，沿单位方向$\vct n$的导数为$\nabla f\cdot\vct n$，不应再乘一个未归一化的位移长度。

散度定义为小体积的单位体积净外通量，旋度的法向分量定义为小回路的单位面积环量：
\[\nabla\cdot\vct A=\lim_{\Delta V\to0}\frac1{\Delta V}\iint_{\partial\Delta V}\vct A\cdot\vct n\dd S,
\quad(\nabla\times\vct A)\cdot\vct n=\lim_{\Delta S\to0}\frac1{\Delta S}\oint\vct A\cdot\dd\vct l.\]
环路方向与法向按右手规则对应。直角坐标中，分别对微小长方体的三对面和三个正交矩形回路相减，得到
\[\nabla\cdot\vct A=\partial_xA_x+\partial_yA_y+\partial_zA_z,\]
\[\nabla\times\vct A=(\partial_yA_z-\partial_zA_y)\vct e_x
+(\partial_zA_x-\partial_xA_z)\vct e_y
+(\partial_xA_y-\partial_yA_x)\vct e_z.\]
通量中的积分对象是曲面，环量中的积分对象是路径。环量为零与通量为零描述不同性质，二者不能互相替代。

\subsection{尺度因子与几何微元}
正交坐标$u_1,u_2,u_3$的尺度因子定义为$h_i=|\partial\vct r/\partial u_i|$，单位基矢为$\vct e_i=(\partial\vct r/\partial u_i)/h_i$。于是
\[\dd\vct r=\sum_i h_i\dd u_i\vct e_i,\qquad
\dd V=h_1h_2h_3\dd u_1\dd u_2\dd u_3.\]
基矢和尺度因子一般都是三个坐标的函数。柱坐标$(s,\phi,z)$有$(h_s,h_\phi,h_z)=(1,s,1)$；球坐标$(r,\theta,\phi)$有$(1,r,r\sin\theta)$，其中$\theta$从正$z$轴量起。

将全微分与线元逐项比较，得到
\[\nabla f=\sum_i\frac1{h_i}\partial_{u_i}f\,\vct e_i.\]
垂直$u_i$方向的两面面积为$(h_1h_2h_3/h_i)\dd u_j\dd u_k$。两面通量相减后除以体积，
\[\nabla\cdot\vct A=\frac1{h_1h_2h_3}
\sum_i\partial_{u_i}\left(\frac{h_1h_2h_3}{h_i}A_i\right).\]
对另外两个方向组成的矩形回路计算环量，同样有
\[(\nabla\times\vct A)_i=
\frac{\partial_{u_j}(h_kA_k)-\partial_{u_k}(h_jA_j)}{h_jh_k},\]
其中$(i,j,k)$按右手次序循环。尺度因子来自相邻曲面的面积和相邻回路的边长，不能只把直角坐标中的字母换掉。

\subsection{球、柱坐标中本讲所用的公式}
球坐标中，
\[\vct E=-\partial_r\varphi\,\vct e_r
-\frac1r\partial_\theta\varphi\,\vct e_\theta
-\frac1{r\sin\theta}\partial_\phi\varphi\,\vct e_\phi.\]
轴对称时最后一项为零，Laplace算子为
\[\nabla^2\varphi=\frac1{r^2}\partial_r(r^2\partial_r\varphi)
+\frac1{r^2\sin\theta}\partial_\theta(\sin\theta\partial_\theta\varphi).\]
柱坐标中，若沿轴平移不变，
\[\nabla^2\varphi=\frac1s\partial_s(s\partial_s\varphi)
+\frac1{s^2}\partial_\phi^2\varphi.\]
轴线上坐标基矢的表示发生退化，不等于场在那里必有物理奇点。需用原矢量场或极限判断。例如均匀场在柱轴处仍有限，而线电荷场的$1/s$才是真实源奇性。


\subsection{把球坐标的散度与旋度写全}
一般矢量$\vct A=A_r\vct e_r+A_\theta\vct e_\theta+A_\phi\vct e_\phi$的球坐标散度是
\[\nabla\cdot\vct A=\frac1{r^2}\partial_r(r^2A_r)
+\frac1{r\sin\theta}\partial_\theta(\sin\theta A_\theta)
+\frac1{r\sin\theta}\partial_\phi A_\phi.
\]
旋度的三个分量为
\[\begin{aligned}
(\nabla\times\vct A)_r&=\frac{\partial_\theta(\sin\theta A_\phi)-\partial_\phi A_\theta}{r\sin\theta},\\
(\nabla\times\vct A)_\theta&=\frac{\partial_\phi A_r}{r\sin\theta}-\frac1r\partial_r(rA_\phi),\\
(\nabla\times\vct A)_\phi&=\frac1r[\partial_r(rA_\theta)-\partial_\theta A_r].
\end{aligned}\]
例如均匀场$E_0\vct e_z$在球坐标中是$E_r=E_0\cos\theta$、$E_\theta=-E_0\sin\theta$。散度中的径向项为$2E_0\cos\theta/r$，角向项为其负值，二者抵消。只对分量作$\partial_rE_r+\partial_\theta E_\theta$，就会得到错误结果。

柱坐标的散度与轴向旋度分别为
\[\nabla\cdot\vct A=\frac1s\partial_s(sA_s)+\frac1s\partial_\phi A_\phi+\partial_zA_z,
\qquad (\nabla\times\vct A)_z=\frac{\partial_s(sA_\phi)-\partial_\phi A_s}{s}.
\]
其他两个旋度分量为$s^{-1}\partial_\phi A_z-\partial_zA_\phi$、$\partial_zA_s-\partial_sA_z$。对不依赖方位角的场，分量的角导数虽为零，基矢随角度变化的贡献仍已包含在这些尺度因子里。

\begin{example}[由有限电荷对得到偶极势]\label{ex:es1-e6}
$+q$与$-q$位于$z=\ell/2$和$z=-\ell/2$。记$\vct p=q\ell\vct e_z$，求精确电势，并在$r\gg\ell$时得到首个非零项。用直接展开和空间平移两种方法解释符号。 见图~\ref{fig:es1-e6-q}。
\end{example}

\begin{figure}[H]\centering
\begin{tikzpicture}[>=Stealth,line cap=round,line join=round]
\draw[axis](-2.7,0)--(2.9,0)node[right]{$z$};\fill(-0.75,0)circle(1.6pt)node[below]{$-q$};\fill(0.75,0)circle(1.6pt)node[below]{$+q$};\fill(0.4,2)circle(1.5pt)node[above right]{$P$};\draw[guide](-0.75,0)--(0.4,2);\draw[guide](0.75,0)--(0.4,2);\draw[axis](-.75,-.6)--(.75,-.6)node[midway,below]{$\vct p$};
\end{tikzpicture}
\caption{两电荷的中点为展开原点，极角从偶极矩正方向量起。}\label{fig:es1-e6-q}
\end{figure}
\begin{solution}
精确距离为$r_\pm=[r^2+\ell^2/4\mp r\ell\cos\theta]^{1/2}$，故
\[\varphi=\ke q\left(\frac1{r_+}-\frac1{r_-}\right).\]
令$a=\ell/(2r)$。只在$|a|\ll1$的远区展开两个平方根倒数，对两个符号分别使用
\[(1+u)^{-1/2}=1-\frac u2+\frac{3u^2}{8}+O(u^3),\quad
u=a^2\mp2a\cos\theta.
\]
保留到一次位移时，$r_\pm^{-1}=r^{-1}[1\pm a\cos\theta+O(a^2)]$，两式相减为$2a\cos\theta/r=\ell\cos\theta/r^2$。中点对称性进一步消去全部偶次项，故得到
\[\boxed{\varphi_{\rm d}(r,\theta)=\ke\frac{p\cos\theta}{r^2}},\qquad
\varphi-\varphi_{\rm d}=O\left(\frac{\ke q\ell^3}{r^4}\right).\]
上述误差为绝对误差。在$\cos\theta=0$附近，偶极首项趋零，不能用它作分母估计相对误差。

图~\ref{fig:es1-e6-s}以$\vct R$表示从源中心指向观察点的矢量。将$1/|\vct R\mp\vct\ell/2|$在$\vct R$处展开，两式相减为$-\vct\ell\cdot\nabla(1/R)$。因为$\nabla(1/R)=-\vct R/R^3$，所以$\varphi_{\rm d}=\ke\vct p\cdot\vct R/R^3$，与上式同号。

若从负电荷所在处作一侧展开，也可把正电荷势与负电荷势看成一次平移之差。此时原点相差$O(\ell)$，但首个偶极项不变，次阶误差的具体形式随原点改变。对精确式比较时，必须明确观察矢量从哪个点量起。
\begin{figure}[H]\centering
\begin{tikzpicture}[>=Stealth,line cap=round,line join=round]

\draw[axis](-.6,0)--(3.3,0)node[right]{$x$};\draw[axis](0,-.8)--(0,2.5)node[above]{$z$};
\fill(0,.3)circle(1.6pt)node[left]{$+q$};\fill(0,-.3)circle(1.6pt)node[left]{$-q$};\fill(2.5,1.8)circle(1.5pt)node[right]{$P$};
\draw[axis](0,0)--(2.5,1.8)node[midway,below]{$\vct R$};\draw[guide](0,.3)--(2.5,1.8) (0,-.3)--(2.5,1.8);\draw[axis](-.8,-.3)--(-.8,.3)node[midway,left]{$\vct\ell$};

\end{tikzpicture}
\caption{从中心到观察点的矢量与两条实际距离分开标记，避免在平移展开中颠倒正负号。}\label{fig:es1-e6-s}
\end{figure}


第三种推导可直接从电场做平移差。点电荷场的方向导数满足
\[E_i\simeq-\ke p_j\partial_j(r_i/r^3)
=\frac{\ke}{r^3}(3n_in_j-\delta_{ij})p_j.\]
使用$\partial_j(r_i/r^3)=\delta_{ij}/r^3-3r_ir_j/r^5$，得到$\vct E=\ke[3(\vct p\cdot\vct n)\vct n-\vct p]/r^3$。再投影到径向和角向，便分别出现$2\cos\theta$与$\sin\theta$，与势法完全相同。
\end{solution}

\begin{example}[偶极场与轴对称场线]\label{ex:es1-e7}
从$\varphi=\ke p\cos\theta/r^2$求理想偶极场。分别由切线方程和球冠通量求场线，说明轴线与赤道上的场方向。 见图~\ref{fig:es1-e7-q}。
\end{example}

\begin{figure}[H]\centering
\begin{tikzpicture}[>=Stealth,line cap=round,line join=round]
\draw[curve](0,0)circle(1.5);\fill(0,0)circle(1.4pt)node[below left]{$O$};\draw[axis](0,-1.9)--(0,2.1)node[above]{$\vct p$};\draw[axis](0,0)--(1.2,.9)node[right]{$r$};\draw(0,.5)arc(90:36.87:.5);\node at(.3,.6){$\theta$};
\end{tikzpicture}
\caption{球坐标极角相对正$z$轴；$\vct p$沿正轴。}\label{fig:es1-e7-q}
\end{figure}
\begin{solution}
球坐标的角向梯度带$1/r$因子。分别对$r$与$\theta$求偏导，
\[\begin{aligned}
E_r&=-\partial_r(\ke p\cos\theta\,r^{-2})
=-\ke p\cos\theta(-2r^{-3}),\\
E_\theta&=-\frac1r\partial_\theta(\ke p\cos\theta\,r^{-2})
=-\frac{\ke p}{r^3}(-\sin\theta).
\end{aligned}\]
故
\[\boxed{E_r=\frac{2\ke p\cos\theta}{r^3},\qquad
E_\theta=\frac{\ke p\sin\theta}{r^3},\qquad E_\phi=0}.\]
沿场线有$\dd r:(r\dd\theta)=E_r:E_\theta$，故
\[\frac{\dd r}r=2\cot\theta\dd\theta,
\qquad \ln r=2\int\frac{\cos\theta}{\sin\theta}\dd\theta+C_0
=2\ln\sin\theta+C_0,
\quad \boxed{r=C\sin^2\theta}.\]
图~\ref{fig:es1-e7-s}画出子午截面中的完整曲线族。$\theta=0,\pi$使角向场为零，轴线需作为单独分支；不能在除以$\sin\theta$后忘记它。

也可积分从正轴到$\theta$的球冠通量：
\[\Psi(r,\theta)=2\pi\int_0^\theta E_r(r,u)r^2\sin u\dd u
=\frac{4\pi\ke p}{r}\left[\frac{\sin^2u}{2}\right]_0^\theta
=\frac{p}{2\epsilon_0r}\sin^2\theta.\]
同一场管的通量不变，故$\sin^2\theta/r$为常数，得到同一场线。此推导既使用轴对称，也使用源外无散度；任意非轴对称场不能照搬球冠积分。

正轴上的场沿正$z$，负轴上的场也沿正$z$，赤道上$E_r=0$而$\vct e_\theta=-\vct e_z$，所以场沿负$z$。在有限电荷对中，近源场线分别终止于真实电荷；理想偶极图把源区压到一点，不能据图把远区公式用于$r\sim\ell$。
\begin{figure}[H]\centering
\begin{tikzpicture}[>=Stealth,line cap=round,line join=round]
\draw[axis](-2.45,0)--(2.6,0)node[right]{$x$};\draw[axis](0,-2.05)--(0,2.2)node[above]{$z$};\draw[curve,domain=1:179,samples=140,smooth]plot ({-1*(sin(\x))^3},{1*(sin(\x))^2*cos(\x)});\draw[->](-0.37222,0.35945)--(-0.52951,0.38471);\draw[->](-0.52951,-0.38471)--(-0.37222,-0.35945);\draw[curve,domain=1:179,samples=140,smooth]plot ({1*(sin(\x))^3},{1*(sin(\x))^2*cos(\x)});\draw[->](0.37222,0.35945)--(0.52951,0.38471);\draw[->](0.52951,-0.38471)--(0.37222,-0.35945);\draw[curve,domain=1:179,samples=140,smooth]plot ({-1.6*(sin(\x))^3},{1.6*(sin(\x))^2*cos(\x)});\draw[->](-0.59556,0.57512)--(-0.84721,0.61554);\draw[->](-0.84721,-0.61554)--(-0.59556,-0.57512);\draw[curve,domain=1:179,samples=140,smooth]plot ({1.6*(sin(\x))^3},{1.6*(sin(\x))^2*cos(\x)});\draw[->](0.59556,0.57512)--(0.84721,0.61554);\draw[->](0.84721,-0.61554)--(0.59556,-0.57512);\draw[curve,domain=1:179,samples=140,smooth]plot ({-2.2*(sin(\x))^3},{2.2*(sin(\x))^2*cos(\x)});\draw[->](-0.81889,0.79079)--(-1.16492,0.84636);\draw[->](-1.16492,-0.84636)--(-0.81889,-0.79079);\draw[curve,domain=1:179,samples=140,smooth]plot ({2.2*(sin(\x))^3},{2.2*(sin(\x))^2*cos(\x)});\draw[->](0.81889,0.79079)--(1.16492,0.84636);\draw[->](1.16492,-0.84636)--(0.81889,-0.79079);\fill(0,0)circle(1.4pt);\draw[->,thick](0,1.05)--(0,1.65)node[right]{$\vct p$};
\end{tikzpicture}
\caption{理想偶极的子午场线；箭头沿电场方向，极角从正$z$轴量起。}\label{fig:es1-e7-s}
\end{figure}


求一条具体曲线时，可用它在赤道的交点半径$r_e$定常数：$C=r_e$。在北半球，沿场方向$\theta$增大，$r$先增大；穿过赤道后$\theta$继续增大而$r$减小，场线回到源区。赤道是该曲线的最远点，不是场线的终点，因为那里角向场最大且非零。若给定通过$(r_0,\theta_0)$的曲线，则$C=r_0/\sin^2\theta_0$，但轴上初值应使用独立的轴线分支。
\end{solution}

\begin{example}[两个反向偶极组成的四极子]\label{ex:es1-e8}
三个点电荷$-q,+2q,-q$分别位于$z=\ell,0,-\ell$。求远区电势及场，分别用精确电荷式和反向偶极的平移说明首项，并给出量级误差。 见图~\ref{fig:es1-e8-q}。
\end{example}

\begin{figure}[H]\centering
\begin{tikzpicture}[>=Stealth,line cap=round,line join=round]
\draw[axis](-2.7,0)--(2.9,0)node[right]{$z$};
\fill(-1.4,0)circle(1.6pt)node[below]{$-q$};
\fill(0,0)circle(1.6pt)node[below]{$+2q$}node[above left]{$O$};
\fill(1.4,0)circle(1.6pt)node[below]{$-q$};
\fill(.6,1.7)circle(1.5pt)node[above right]{$P$};
\draw[guide](-1.4,0)--(.6,1.7);\draw[guide](1.4,0)--(.6,1.7);
\draw[axis](0,0)--(.6,1.7)node[midway,left,fill=white,inner sep=1pt]{$r$};
\draw(.42,0)arc(0:70.56:.42);\node at(.48,.34){$\theta$};
\draw[<->](-1.4,-.7)--(0,-.7)node[midway,below]{$\ell$};
\draw[<->](0,-.7)--(1.4,-.7)node[midway,below]{$\ell$};
\end{tikzpicture}
\caption{总电荷与偶极矩均为零，电荷分布关于中点反演对称。}\label{fig:es1-e8-q}
\end{figure}
\begin{solution}
精确电势为
\[\varphi=\ke q\left[\frac2r-
\frac1{\sqrt{r^2+\ell^2-2r\ell\cos\theta}}
-\frac1{\sqrt{r^2+\ell^2+2r\ell\cos\theta}}\right].\]
两个分母相加时奇次项抵消，常数项再与$2/r$抵消。具体地，令$t=\ell/r$、$c=\cos\theta$，则
\[(1\mp2tc+t^2)^{-1/2}
=1\pm tc+\frac{t^2}{2}(3c^2-1)+O(t^3).
\]
两式相加后为$2+t^2(3c^2-1)+O(t^4)$，再乘$1/r$并从$2/r$中减去。保留二阶位移项，
\[\boxed{\varphi_{\rm q}=-\frac{\ke q\ell^2}{r^3}(3\cos^2\theta-1)},\qquad
\varphi-\varphi_{\rm q}=O(\ke q\ell^4/r^5).\]
记$A=q\ell^2$。径向求导只作用于$r^{-3}$，角向求导则作用于$3\cos^2\theta-1$，故
\[E_r=\ke A(3\cos^2\theta-1)(-3r^{-4}),\qquad
E_\theta=\frac{\ke A}{r^4}\,3(2\cos\theta)(-\sin\theta).\]
于是
\[\boxed{E_r=-\frac{3\ke A}{r^4}(3\cos^2\theta-1),\qquad
E_\theta=-\frac{6\ke A}{r^4}\sin\theta\cos\theta}.\]
若改写成二倍角，$E_r=-3\ke A(3\cos2\theta+1)/(2r^4)$，$E_\theta=-3\ke A\sin2\theta/r^4$。角度改写中的$1/2$与径向求导中的$3$来自不同步骤，应分别保留。

图~\ref{fig:es1-e8-s}中的$+2q$可分为两个$+q$：其中一个与$z=\ell$的$-q$组成向下偶极，另一个与$z=-\ell$的$-q$组成向上偶极。两矩相反、中心相隔$\ell$；向上偶极在下，向下偶极在上。记单个偶极的矩大小为$p=q\ell$，向上偶极势为$\varphi_{\rm d}=\ke p\cos\theta/r^2$。两者的平移差首项为$\ell\partial_z\varphi_{\rm d}$，使用
\[\partial_z=\cos\theta\partial_r-\frac{\sin\theta}{r}\partial_\theta,\]
得到
\[\begin{aligned}
\partial_z\varphi_{\rm d}
&=\cos\theta\frac{-2\ke p\cos\theta}{r^3}
-\frac{\sin\theta}{r}\frac{-\ke p\sin\theta}{r^2}\\
&=\frac{\ke p}{r^3}(\sin^2\theta-2\cos^2\theta)
=-\frac{\ke p}{r^3}(3\cos^2\theta-1).
\end{aligned}\]
再乘$\ell$便恢复四极势。平移差多出一个$\ell/r$，使衰减由$r^{-2}$变为$r^{-3}$；其负号由两偶极的上下次序确定，整个三电荷体系的净偶极矩仍为零。
\begin{figure}[H]\centering
\begin{tikzpicture}[>=Stealth,line cap=round,line join=round]

\draw[axis](0,-2)--(0,2.2)node[above]{$z$};\fill(0,1.4)circle(1.7pt)node[right]{$-q$};\fill(0,0)circle(1.7pt)node[right]{$+2q$};\fill(0,-1.4)circle(1.7pt)node[right]{$-q$};
\draw[axis](-.55,1.25)--(-.55,.15)node[midway,left]{$-\vct p$};\draw[axis](-.55,-1.25)--(-.55,-.15)node[midway,left]{$\vct p$};\node[right,align=left]at(1.6,.5){$Q=0$\\[5pt]$\vct p_{\rm total}=0$\\[5pt]$A=q\ell^2$};

\end{tikzpicture}
\caption{两个反向偶极的一阶平移差产生四极首项，不能把单个偶极矩误记成整个体系的偶极矩。}\label{fig:es1-e8-s}
\end{figure}


\end{solution}

\begin{example}[四极场线的分支]\label{ex:es1-e9}
对例~\ref{ex:es1-e8}的理想四极场，求场线方程。说明南北半球的积分常数怎样选择，并用球冠通量验证。 见图~\ref{fig:es1-e9-q}。
\end{example}

\begin{figure}[H]\centering
\begin{tikzpicture}[>=Stealth,line cap=round,line join=round]
\draw[axis](-2.3,0)--(2.45,0)node[right]{$x$};\draw[axis](0,-2.3)--(0,2.45)node[above]{$z$};\draw[curve]plot coordinates{(-0.000001,0.001000) (-0.000278,0.016684) (-0.001048,0.032345) (-0.002307,0.047962) (-0.004055,0.063515) (-0.006287,0.078981) (-0.009002,0.094341) (-0.012192,0.109573) (-0.015854,0.124657) (-0.019981,0.139573) (-0.024565,0.154301) (-0.029599,0.168820) (-0.035075,0.183112) (-0.040981,0.197157) (-0.047309,0.210937) (-0.054048,0.224433) (-0.061185,0.237628) (-0.068708,0.250505) (-0.076604,0.263045) (-0.084860,0.275234) (-0.093460,0.287056) (-0.102391,0.298494) (-0.111635,0.309535) (-0.121177,0.320164) (-0.131001,0.330368) (-0.141088,0.340135) (-0.151421,0.349452) (-0.161981,0.358309) (-0.172750,0.366694) (-0.183709,0.374598) (-0.194837,0.382011) (-0.206115,0.388926) (-0.217522,0.395334) (-0.229038,0.401230) (-0.240642,0.406606) (-0.252312,0.411459) (-0.264027,0.415783) (-0.275766,0.419576) (-0.287506,0.422834) (-0.299225,0.425555) (-0.310902,0.427740) (-0.322514,0.429388) (-0.334039,0.430499) (-0.345455,0.431075) (-0.356740,0.431119) (-0.367870,0.430634) (-0.378825,0.429624) (-0.389582,0.428094) (-0.400120,0.426050) (-0.410416,0.423499) (-0.420448,0.420448) (-0.430196,0.416906) (-0.439639,0.412881) (-0.448754,0.408384) (-0.457523,0.403426) (-0.465923,0.398017) (-0.473936,0.392170) (-0.481541,0.385898) (-0.488719,0.379214) (-0.495450,0.372134) (-0.501716,0.364672) (-0.507499,0.356843) (-0.512780,0.348665) (-0.517541,0.340154) (-0.521766,0.331328) (-0.525437,0.322205) (-0.528537,0.312804) (-0.531050,0.303145) (-0.532960,0.293247) (-0.534251,0.283131) (-0.534907,0.272818) (-0.534913,0.262330) (-0.534252,0.251687) (-0.532910,0.240913) (-0.530872,0.230031) (-0.528120,0.219064) (-0.524639,0.208035) (-0.520412,0.196969) (-0.515422,0.185890) (-0.509651,0.174823) (-0.503078,0.163794) (-0.495681,0.152828) (-0.487438,0.141952) (-0.478321,0.131193) (-0.468301,0.120579) (-0.457343,0.110137) (-0.445405,0.099896) (-0.432439,0.089888) (-0.418386,0.080141) (-0.403175,0.070689) (-0.386714,0.061567) (-0.368888,0.052809) (-0.349545,0.044456) (-0.328482,0.036551) (-0.305415,0.029141) (-0.279928,0.022284) (-0.251381,0.016048) (-0.218687,0.010519) (-0.179729,0.005821) (-0.129143,0.002155) (-0.031623,0.000032)};\draw[curve]plot coordinates{(-0.031623,-0.000032) (-0.129143,-0.002155) (-0.179729,-0.005821) (-0.218687,-0.010519) (-0.251381,-0.016048) (-0.279928,-0.022284) (-0.305415,-0.029141) (-0.328482,-0.036551) (-0.349545,-0.044456) (-0.368888,-0.052809) (-0.386714,-0.061567) (-0.403175,-0.070689) (-0.418386,-0.080141) (-0.432439,-0.089888) (-0.445405,-0.099896) (-0.457343,-0.110137) (-0.468301,-0.120579) (-0.478321,-0.131193) (-0.487438,-0.141952) (-0.495681,-0.152828) (-0.503078,-0.163794) (-0.509651,-0.174823) (-0.515422,-0.185890) (-0.520412,-0.196969) (-0.524639,-0.208035) (-0.528120,-0.219064) (-0.530872,-0.230031) (-0.532910,-0.240913) (-0.534252,-0.251687) (-0.534913,-0.262330) (-0.534907,-0.272818) (-0.534251,-0.283131) (-0.532960,-0.293247) (-0.531050,-0.303145) (-0.528537,-0.312804) (-0.525437,-0.322205) (-0.521766,-0.331328) (-0.517541,-0.340154) (-0.512780,-0.348665) (-0.507499,-0.356843) (-0.501716,-0.364672) (-0.495450,-0.372134) (-0.488719,-0.379214) (-0.481541,-0.385898) (-0.473936,-0.392170) (-0.465923,-0.398017) (-0.457523,-0.403426) (-0.448754,-0.408384) (-0.439639,-0.412881) (-0.430196,-0.416906) (-0.420448,-0.420448) (-0.410416,-0.423499) (-0.400120,-0.426050) (-0.389582,-0.428094) (-0.378825,-0.429624) (-0.367870,-0.430634) (-0.356740,-0.431119) (-0.345455,-0.431075) (-0.334039,-0.430499) (-0.322514,-0.429388) (-0.310902,-0.427740) (-0.299225,-0.425555) (-0.287506,-0.422834) (-0.275766,-0.419576) (-0.264027,-0.415783) (-0.252312,-0.411459) (-0.240642,-0.406606) (-0.229038,-0.401230) (-0.217522,-0.395334) (-0.206115,-0.388926) (-0.194837,-0.382011) (-0.183709,-0.374598) (-0.172750,-0.366694) (-0.161981,-0.358309) (-0.151421,-0.349452) (-0.141088,-0.340135) (-0.131001,-0.330368) (-0.121177,-0.320164) (-0.111635,-0.309535) (-0.102391,-0.298494) (-0.093460,-0.287056) (-0.084860,-0.275234) (-0.076604,-0.263045) (-0.068708,-0.250505) (-0.061185,-0.237628) (-0.054048,-0.224433) (-0.047309,-0.210937) (-0.040981,-0.197157) (-0.035075,-0.183112) (-0.029599,-0.168820) (-0.024565,-0.154301) (-0.019981,-0.139573) (-0.015854,-0.124657) (-0.012192,-0.109573) (-0.009002,-0.094341) (-0.006287,-0.078981) (-0.004055,-0.063515) (-0.002307,-0.047962) (-0.001048,-0.032345) (-0.000278,-0.016684) (-0.000001,-0.001000)};\draw[curve]plot coordinates{(0.000001,0.001000) (0.000278,0.016684) (0.001048,0.032345) (0.002307,0.047962) (0.004055,0.063515) (0.006287,0.078981) (0.009002,0.094341) (0.012192,0.109573) (0.015854,0.124657) (0.019981,0.139573) (0.024565,0.154301) (0.029599,0.168820) (0.035075,0.183112) (0.040981,0.197157) (0.047309,0.210937) (0.054048,0.224433) (0.061185,0.237628) (0.068708,0.250505) (0.076604,0.263045) (0.084860,0.275234) (0.093460,0.287056) (0.102391,0.298494) (0.111635,0.309535) (0.121177,0.320164) (0.131001,0.330368) (0.141088,0.340135) (0.151421,0.349452) (0.161981,0.358309) (0.172750,0.366694) (0.183709,0.374598) (0.194837,0.382011) (0.206115,0.388926) (0.217522,0.395334) (0.229038,0.401230) (0.240642,0.406606) (0.252312,0.411459) (0.264027,0.415783) (0.275766,0.419576) (0.287506,0.422834) (0.299225,0.425555) (0.310902,0.427740) (0.322514,0.429388) (0.334039,0.430499) (0.345455,0.431075) (0.356740,0.431119) (0.367870,0.430634) (0.378825,0.429624) (0.389582,0.428094) (0.400120,0.426050) (0.410416,0.423499) (0.420448,0.420448) (0.430196,0.416906) (0.439639,0.412881) (0.448754,0.408384) (0.457523,0.403426) (0.465923,0.398017) (0.473936,0.392170) (0.481541,0.385898) (0.488719,0.379214) (0.495450,0.372134) (0.501716,0.364672) (0.507499,0.356843) (0.512780,0.348665) (0.517541,0.340154) (0.521766,0.331328) (0.525437,0.322205) (0.528537,0.312804) (0.531050,0.303145) (0.532960,0.293247) (0.534251,0.283131) (0.534907,0.272818) (0.534913,0.262330) (0.534252,0.251687) (0.532910,0.240913) (0.530872,0.230031) (0.528120,0.219064) (0.524639,0.208035) (0.520412,0.196969) (0.515422,0.185890) (0.509651,0.174823) (0.503078,0.163794) (0.495681,0.152828) (0.487438,0.141952) (0.478321,0.131193) (0.468301,0.120579) (0.457343,0.110137) (0.445405,0.099896) (0.432439,0.089888) (0.418386,0.080141) (0.403175,0.070689) (0.386714,0.061567) (0.368888,0.052809) (0.349545,0.044456) (0.328482,0.036551) (0.305415,0.029141) (0.279928,0.022284) (0.251381,0.016048) (0.218687,0.010519) (0.179729,0.005821) (0.129143,0.002155) (0.031623,0.000032)};\draw[curve]plot coordinates{(0.031623,-0.000032) (0.129143,-0.002155) (0.179729,-0.005821) (0.218687,-0.010519) (0.251381,-0.016048) (0.279928,-0.022284) (0.305415,-0.029141) (0.328482,-0.036551) (0.349545,-0.044456) (0.368888,-0.052809) (0.386714,-0.061567) (0.403175,-0.070689) (0.418386,-0.080141) (0.432439,-0.089888) (0.445405,-0.099896) (0.457343,-0.110137) (0.468301,-0.120579) (0.478321,-0.131193) (0.487438,-0.141952) (0.495681,-0.152828) (0.503078,-0.163794) (0.509651,-0.174823) (0.515422,-0.185890) (0.520412,-0.196969) (0.524639,-0.208035) (0.528120,-0.219064) (0.530872,-0.230031) (0.532910,-0.240913) (0.534252,-0.251687) (0.534913,-0.262330) (0.534907,-0.272818) (0.534251,-0.283131) (0.532960,-0.293247) (0.531050,-0.303145) (0.528537,-0.312804) (0.525437,-0.322205) (0.521766,-0.331328) (0.517541,-0.340154) (0.512780,-0.348665) (0.507499,-0.356843) (0.501716,-0.364672) (0.495450,-0.372134) (0.488719,-0.379214) (0.481541,-0.385898) (0.473936,-0.392170) (0.465923,-0.398017) (0.457523,-0.403426) (0.448754,-0.408384) (0.439639,-0.412881) (0.430196,-0.416906) (0.420448,-0.420448) (0.410416,-0.423499) (0.400120,-0.426050) (0.389582,-0.428094) (0.378825,-0.429624) (0.367870,-0.430634) (0.356740,-0.431119) (0.345455,-0.431075) (0.334039,-0.430499) (0.322514,-0.429388) (0.310902,-0.427740) (0.299225,-0.425555) (0.287506,-0.422834) (0.275766,-0.419576) (0.264027,-0.415783) (0.252312,-0.411459) (0.240642,-0.406606) (0.229038,-0.401230) (0.217522,-0.395334) (0.206115,-0.388926) (0.194837,-0.382011) (0.183709,-0.374598) (0.172750,-0.366694) (0.161981,-0.358309) (0.151421,-0.349452) (0.141088,-0.340135) (0.131001,-0.330368) (0.121177,-0.320164) (0.111635,-0.309535) (0.102391,-0.298494) (0.093460,-0.287056) (0.084860,-0.275234) (0.076604,-0.263045) (0.068708,-0.250505) (0.061185,-0.237628) (0.054048,-0.224433) (0.047309,-0.210937) (0.040981,-0.197157) (0.035075,-0.183112) (0.029599,-0.168820) (0.024565,-0.154301) (0.019981,-0.139573) (0.015854,-0.124657) (0.012192,-0.109573) (0.009002,-0.094341) (0.006287,-0.078981) (0.004055,-0.063515) (0.002307,-0.047962) (0.001048,-0.032345) (0.000278,-0.016684) (0.000001,-0.001000)};\draw[curve]plot coordinates{(-0.000002,0.001600) (-0.000446,0.026694) (-0.001676,0.051752) (-0.003691,0.076739) (-0.006487,0.101623) (-0.010060,0.126370) (-0.014403,0.150945) (-0.019508,0.175317) (-0.025367,0.199452) (-0.031970,0.223317) (-0.039305,0.246881) (-0.047359,0.270112) (-0.056119,0.292979) (-0.065570,0.315451) (-0.075695,0.337499) (-0.086477,0.359093) (-0.097896,0.380205) (-0.109933,0.400807) (-0.122567,0.420873) (-0.135776,0.440375) (-0.149537,0.459289) (-0.163825,0.477590) (-0.178616,0.495255) (-0.193884,0.512262) (-0.209601,0.528589) (-0.225741,0.544216) (-0.242273,0.559124) (-0.259170,0.573294) (-0.276400,0.586710) (-0.293934,0.599356) (-0.311739,0.611218) (-0.329784,0.622281) (-0.348036,0.632535) (-0.366461,0.641968) (-0.385027,0.650570) (-0.403700,0.658334) (-0.422444,0.665253) (-0.441225,0.671321) (-0.460009,0.676534) (-0.478760,0.680889) (-0.497443,0.684384) (-0.516023,0.687020) (-0.534463,0.688798) (-0.552728,0.689720) (-0.570783,0.689790) (-0.588592,0.689014) (-0.606120,0.687398) (-0.623332,0.684950) (-0.640192,0.681680) (-0.656665,0.677598) (-0.672717,0.672717) (-0.688314,0.667050) (-0.703422,0.660610) (-0.718007,0.653415) (-0.732036,0.645481) (-0.745477,0.636827) (-0.758298,0.627472) (-0.770465,0.617436) (-0.781950,0.606743) (-0.792720,0.595414) (-0.802746,0.583474) (-0.811998,0.570949) (-0.820448,0.557864) (-0.828066,0.544246) (-0.834826,0.530125) (-0.840699,0.515528) (-0.845660,0.500487) (-0.849681,0.485032) (-0.852737,0.469196) (-0.854802,0.453010) (-0.855851,0.436509) (-0.855860,0.419727) (-0.854804,0.402700) (-0.852657,0.385462) (-0.849394,0.368050) (-0.844992,0.350502) (-0.839422,0.332856) (-0.832660,0.315150) (-0.824676,0.297423) (-0.815441,0.279716) (-0.804924,0.262070) (-0.793090,0.244525) (-0.779901,0.227124) (-0.765314,0.209909) (-0.749282,0.192926) (-0.731748,0.176219) (-0.712647,0.159834) (-0.691902,0.143820) (-0.669418,0.128226) (-0.645080,0.113103) (-0.618742,0.098507) (-0.590221,0.084495) (-0.559272,0.071130) (-0.525572,0.058481) (-0.488664,0.046626) (-0.447886,0.035655) (-0.402210,0.025676) (-0.349899,0.016830) (-0.287566,0.009313) (-0.206629,0.003449) (-0.050596,0.000051)};\draw[curve]plot coordinates{(-0.050596,-0.000051) (-0.206629,-0.003449) (-0.287566,-0.009313) (-0.349899,-0.016830) (-0.402210,-0.025676) (-0.447886,-0.035655) (-0.488664,-0.046626) (-0.525572,-0.058481) (-0.559272,-0.071130) (-0.590221,-0.084495) (-0.618742,-0.098507) (-0.645080,-0.113103) (-0.669418,-0.128226) (-0.691902,-0.143820) (-0.712647,-0.159834) (-0.731748,-0.176219) (-0.749282,-0.192926) (-0.765314,-0.209909) (-0.779901,-0.227124) (-0.793090,-0.244525) (-0.804924,-0.262070) (-0.815441,-0.279716) (-0.824676,-0.297423) (-0.832660,-0.315150) (-0.839422,-0.332856) (-0.844992,-0.350502) (-0.849394,-0.368050) (-0.852657,-0.385462) (-0.854804,-0.402700) (-0.855860,-0.419727) (-0.855851,-0.436509) (-0.854802,-0.453010) (-0.852737,-0.469196) (-0.849681,-0.485032) (-0.845660,-0.500487) (-0.840699,-0.515528) (-0.834826,-0.530125) (-0.828066,-0.544246) (-0.820448,-0.557864) (-0.811998,-0.570949) (-0.802746,-0.583474) (-0.792720,-0.595414) (-0.781950,-0.606743) (-0.770465,-0.617436) (-0.758298,-0.627472) (-0.745477,-0.636827) (-0.732036,-0.645481) (-0.718007,-0.653415) (-0.703422,-0.660610) (-0.688314,-0.667050) (-0.672717,-0.672717) (-0.656665,-0.677598) (-0.640192,-0.681680) (-0.623332,-0.684950) (-0.606120,-0.687398) (-0.588592,-0.689014) (-0.570783,-0.689790) (-0.552728,-0.689720) (-0.534463,-0.688798) (-0.516023,-0.687020) (-0.497443,-0.684384) (-0.478760,-0.680889) (-0.460009,-0.676534) (-0.441225,-0.671321) (-0.422444,-0.665253) (-0.403700,-0.658334) (-0.385027,-0.650570) (-0.366461,-0.641968) (-0.348036,-0.632535) (-0.329784,-0.622281) (-0.311739,-0.611218) (-0.293934,-0.599356) (-0.276400,-0.586710) (-0.259170,-0.573294) (-0.242273,-0.559124) (-0.225741,-0.544216) (-0.209601,-0.528589) (-0.193884,-0.512262) (-0.178616,-0.495255) (-0.163825,-0.477590) (-0.149537,-0.459289) (-0.135776,-0.440375) (-0.122567,-0.420873) (-0.109933,-0.400807) (-0.097896,-0.380205) (-0.086477,-0.359093) (-0.075695,-0.337499) (-0.065570,-0.315451) (-0.056119,-0.292979) (-0.047359,-0.270112) (-0.039305,-0.246881) (-0.031970,-0.223317) (-0.025367,-0.199452) (-0.019508,-0.175317) (-0.014403,-0.150945) (-0.010060,-0.126370) (-0.006487,-0.101623) (-0.003691,-0.076739) (-0.001676,-0.051752) (-0.000446,-0.026694) (-0.000002,-0.001600)};\draw[curve]plot coordinates{(0.000002,0.001600) (0.000446,0.026694) (0.001676,0.051752) (0.003691,0.076739) (0.006487,0.101623) (0.010060,0.126370) (0.014403,0.150945) (0.019508,0.175317) (0.025367,0.199452) (0.031970,0.223317) (0.039305,0.246881) (0.047359,0.270112) (0.056119,0.292979) (0.065570,0.315451) (0.075695,0.337499) (0.086477,0.359093) (0.097896,0.380205) (0.109933,0.400807) (0.122567,0.420873) (0.135776,0.440375) (0.149537,0.459289) (0.163825,0.477590) (0.178616,0.495255) (0.193884,0.512262) (0.209601,0.528589) (0.225741,0.544216) (0.242273,0.559124) (0.259170,0.573294) (0.276400,0.586710) (0.293934,0.599356) (0.311739,0.611218) (0.329784,0.622281) (0.348036,0.632535) (0.366461,0.641968) (0.385027,0.650570) (0.403700,0.658334) (0.422444,0.665253) (0.441225,0.671321) (0.460009,0.676534) (0.478760,0.680889) (0.497443,0.684384) (0.516023,0.687020) (0.534463,0.688798) (0.552728,0.689720) (0.570783,0.689790) (0.588592,0.689014) (0.606120,0.687398) (0.623332,0.684950) (0.640192,0.681680) (0.656665,0.677598) (0.672717,0.672717) (0.688314,0.667050) (0.703422,0.660610) (0.718007,0.653415) (0.732036,0.645481) (0.745477,0.636827) (0.758298,0.627472) (0.770465,0.617436) (0.781950,0.606743) (0.792720,0.595414) (0.802746,0.583474) (0.811998,0.570949) (0.820448,0.557864) (0.828066,0.544246) (0.834826,0.530125) (0.840699,0.515528) (0.845660,0.500487) (0.849681,0.485032) (0.852737,0.469196) (0.854802,0.453010) (0.855851,0.436509) (0.855860,0.419727) (0.854804,0.402700) (0.852657,0.385462) (0.849394,0.368050) (0.844992,0.350502) (0.839422,0.332856) (0.832660,0.315150) (0.824676,0.297423) (0.815441,0.279716) (0.804924,0.262070) (0.793090,0.244525) (0.779901,0.227124) (0.765314,0.209909) (0.749282,0.192926) (0.731748,0.176219) (0.712647,0.159834) (0.691902,0.143820) (0.669418,0.128226) (0.645080,0.113103) (0.618742,0.098507) (0.590221,0.084495) (0.559272,0.071130) (0.525572,0.058481) (0.488664,0.046626) (0.447886,0.035655) (0.402210,0.025676) (0.349899,0.016830) (0.287566,0.009313) (0.206629,0.003449) (0.050596,0.000051)};\draw[curve]plot coordinates{(0.050596,-0.000051) (0.206629,-0.003449) (0.287566,-0.009313) (0.349899,-0.016830) (0.402210,-0.025676) (0.447886,-0.035655) (0.488664,-0.046626) (0.525572,-0.058481) (0.559272,-0.071130) (0.590221,-0.084495) (0.618742,-0.098507) (0.645080,-0.113103) (0.669418,-0.128226) (0.691902,-0.143820) (0.712647,-0.159834) (0.731748,-0.176219) (0.749282,-0.192926) (0.765314,-0.209909) (0.779901,-0.227124) (0.793090,-0.244525) (0.804924,-0.262070) (0.815441,-0.279716) (0.824676,-0.297423) (0.832660,-0.315150) (0.839422,-0.332856) (0.844992,-0.350502) (0.849394,-0.368050) (0.852657,-0.385462) (0.854804,-0.402700) (0.855860,-0.419727) (0.855851,-0.436509) (0.854802,-0.453010) (0.852737,-0.469196) (0.849681,-0.485032) (0.845660,-0.500487) (0.840699,-0.515528) (0.834826,-0.530125) (0.828066,-0.544246) (0.820448,-0.557864) (0.811998,-0.570949) (0.802746,-0.583474) (0.792720,-0.595414) (0.781950,-0.606743) (0.770465,-0.617436) (0.758298,-0.627472) (0.745477,-0.636827) (0.732036,-0.645481) (0.718007,-0.653415) (0.703422,-0.660610) (0.688314,-0.667050) (0.672717,-0.672717) (0.656665,-0.677598) (0.640192,-0.681680) (0.623332,-0.684950) (0.606120,-0.687398) (0.588592,-0.689014) (0.570783,-0.689790) (0.552728,-0.689720) (0.534463,-0.688798) (0.516023,-0.687020) (0.497443,-0.684384) (0.478760,-0.680889) (0.460009,-0.676534) (0.441225,-0.671321) (0.422444,-0.665253) (0.403700,-0.658334) (0.385027,-0.650570) (0.366461,-0.641968) (0.348036,-0.632535) (0.329784,-0.622281) (0.311739,-0.611218) (0.293934,-0.599356) (0.276400,-0.586710) (0.259170,-0.573294) (0.242273,-0.559124) (0.225741,-0.544216) (0.209601,-0.528589) (0.193884,-0.512262) (0.178616,-0.495255) (0.163825,-0.477590) (0.149537,-0.459289) (0.135776,-0.440375) (0.122567,-0.420873) (0.109933,-0.400807) (0.097896,-0.380205) (0.086477,-0.359093) (0.075695,-0.337499) (0.065570,-0.315451) (0.056119,-0.292979) (0.047359,-0.270112) (0.039305,-0.246881) (0.031970,-0.223317) (0.025367,-0.199452) (0.019508,-0.175317) (0.014403,-0.150945) (0.010060,-0.126370) (0.006487,-0.101623) (0.003691,-0.076739) (0.001676,-0.051752) (0.000446,-0.026694) (0.000002,-0.001600)};\draw[curve]plot coordinates{(-0.000002,0.002200) (-0.000613,0.036704) (-0.002305,0.071159) (-0.005075,0.105517) (-0.008920,0.139732) (-0.013832,0.173759) (-0.019804,0.207550) (-0.026823,0.241061) (-0.034879,0.274246) (-0.043958,0.307061) (-0.054044,0.339461) (-0.065119,0.371404) (-0.077164,0.402846) (-0.090159,0.433745) (-0.104081,0.464061) (-0.118905,0.493753) (-0.134607,0.522782) (-0.151158,0.551110) (-0.168530,0.578700) (-0.186692,0.605515) (-0.205613,0.631522) (-0.225260,0.656686) (-0.245597,0.680976) (-0.266590,0.704361) (-0.288202,0.726810) (-0.310393,0.748297) (-0.333126,0.768795) (-0.356359,0.788280) (-0.380050,0.806727) (-0.404159,0.824115) (-0.428641,0.840425) (-0.453453,0.855637) (-0.478549,0.869735) (-0.503884,0.882706) (-0.529412,0.894534) (-0.555087,0.905210) (-0.580860,0.914723) (-0.606685,0.923066) (-0.632513,0.930234) (-0.658296,0.936222) (-0.683985,0.941028) (-0.709531,0.944653) (-0.734886,0.947097) (-0.760001,0.948365) (-0.784827,0.948461) (-0.809315,0.947394) (-0.833416,0.945172) (-0.857081,0.941806) (-0.880263,0.937310) (-0.902914,0.931698) (-0.924986,0.924986) (-0.946432,0.917193) (-0.967205,0.908339) (-0.987260,0.898446) (-1.006550,0.887536) (-1.025031,0.875637) (-1.042659,0.862773) (-1.059390,0.848975) (-1.075181,0.834271) (-1.089990,0.818694) (-1.103776,0.802277) (-1.116498,0.785055) (-1.128116,0.767063) (-1.138591,0.748338) (-1.147886,0.728921) (-1.155961,0.708851) (-1.162782,0.688170) (-1.168311,0.666919) (-1.172513,0.645144) (-1.175353,0.622889) (-1.176796,0.600200) (-1.176808,0.577125) (-1.175355,0.553712) (-1.172403,0.530010) (-1.167917,0.506069) (-1.161864,0.481940) (-1.154206,0.457677) (-1.144907,0.433331) (-1.133929,0.408957) (-1.121232,0.384610) (-1.106771,0.360346) (-1.090499,0.336221) (-1.072364,0.312295) (-1.052307,0.288625) (-1.030263,0.265273) (-1.006154,0.242301) (-0.979890,0.219772) (-0.951365,0.197753) (-0.920450,0.176310) (-0.886984,0.155517) (-0.850770,0.135447) (-0.811553,0.116180) (-0.769000,0.097804) (-0.722661,0.080412) (-0.671913,0.064111) (-0.615843,0.049026) (-0.553039,0.035305) (-0.481111,0.023142) (-0.395404,0.012806) (-0.284114,0.004742) (-0.069570,0.000070)};\draw[curve]plot coordinates{(-0.069570,-0.000070) (-0.284114,-0.004742) (-0.395404,-0.012806) (-0.481111,-0.023142) (-0.553039,-0.035305) (-0.615843,-0.049026) (-0.671913,-0.064111) (-0.722661,-0.080412) (-0.769000,-0.097804) (-0.811553,-0.116180) (-0.850770,-0.135447) (-0.886984,-0.155517) (-0.920450,-0.176310) (-0.951365,-0.197753) (-0.979890,-0.219772) (-1.006154,-0.242301) (-1.030263,-0.265273) (-1.052307,-0.288625) (-1.072364,-0.312295) (-1.090499,-0.336221) (-1.106771,-0.360346) (-1.121232,-0.384610) (-1.133929,-0.408957) (-1.144907,-0.433331) (-1.154206,-0.457677) (-1.161864,-0.481940) (-1.167917,-0.506069) (-1.172403,-0.530010) (-1.175355,-0.553712) (-1.176808,-0.577125) (-1.176796,-0.600200) (-1.175353,-0.622889) (-1.172513,-0.645144) (-1.168311,-0.666919) (-1.162782,-0.688170) (-1.155961,-0.708851) (-1.147886,-0.728921) (-1.138591,-0.748338) (-1.128116,-0.767063) (-1.116498,-0.785055) (-1.103776,-0.802277) (-1.089990,-0.818694) (-1.075181,-0.834271) (-1.059390,-0.848975) (-1.042659,-0.862773) (-1.025031,-0.875637) (-1.006550,-0.887536) (-0.987260,-0.898446) (-0.967205,-0.908339) (-0.946432,-0.917193) (-0.924986,-0.924986) (-0.902914,-0.931698) (-0.880263,-0.937310) (-0.857081,-0.941806) (-0.833416,-0.945172) (-0.809315,-0.947394) (-0.784827,-0.948461) (-0.760001,-0.948365) (-0.734886,-0.947097) (-0.709531,-0.944653) (-0.683985,-0.941028) (-0.658296,-0.936222) (-0.632513,-0.930234) (-0.606685,-0.923066) (-0.580860,-0.914723) (-0.555087,-0.905210) (-0.529412,-0.894534) (-0.503884,-0.882706) (-0.478549,-0.869735) (-0.453453,-0.855637) (-0.428641,-0.840425) (-0.404159,-0.824115) (-0.380050,-0.806727) (-0.356359,-0.788280) (-0.333126,-0.768795) (-0.310393,-0.748297) (-0.288202,-0.726810) (-0.266590,-0.704361) (-0.245597,-0.680976) (-0.225260,-0.656686) (-0.205613,-0.631522) (-0.186692,-0.605515) (-0.168530,-0.578700) (-0.151158,-0.551110) (-0.134607,-0.522782) (-0.118905,-0.493753) (-0.104081,-0.464061) (-0.090159,-0.433745) (-0.077164,-0.402846) (-0.065119,-0.371404) (-0.054044,-0.339461) (-0.043958,-0.307061) (-0.034879,-0.274246) (-0.026823,-0.241061) (-0.019804,-0.207550) (-0.013832,-0.173759) (-0.008920,-0.139732) (-0.005075,-0.105517) (-0.002305,-0.071159) (-0.000613,-0.036704) (-0.000002,-0.002200)};\draw[curve]plot coordinates{(0.000002,0.002200) (0.000613,0.036704) (0.002305,0.071159) (0.005075,0.105517) (0.008920,0.139732) (0.013832,0.173759) (0.019804,0.207550) (0.026823,0.241061) (0.034879,0.274246) (0.043958,0.307061) (0.054044,0.339461) (0.065119,0.371404) (0.077164,0.402846) (0.090159,0.433745) (0.104081,0.464061) (0.118905,0.493753) (0.134607,0.522782) (0.151158,0.551110) (0.168530,0.578700) (0.186692,0.605515) (0.205613,0.631522) (0.225260,0.656686) (0.245597,0.680976) (0.266590,0.704361) (0.288202,0.726810) (0.310393,0.748297) (0.333126,0.768795) (0.356359,0.788280) (0.380050,0.806727) (0.404159,0.824115) (0.428641,0.840425) (0.453453,0.855637) (0.478549,0.869735) (0.503884,0.882706) (0.529412,0.894534) (0.555087,0.905210) (0.580860,0.914723) (0.606685,0.923066) (0.632513,0.930234) (0.658296,0.936222) (0.683985,0.941028) (0.709531,0.944653) (0.734886,0.947097) (0.760001,0.948365) (0.784827,0.948461) (0.809315,0.947394) (0.833416,0.945172) (0.857081,0.941806) (0.880263,0.937310) (0.902914,0.931698) (0.924986,0.924986) (0.946432,0.917193) (0.967205,0.908339) (0.987260,0.898446) (1.006550,0.887536) (1.025031,0.875637) (1.042659,0.862773) (1.059390,0.848975) (1.075181,0.834271) (1.089990,0.818694) (1.103776,0.802277) (1.116498,0.785055) (1.128116,0.767063) (1.138591,0.748338) (1.147886,0.728921) (1.155961,0.708851) (1.162782,0.688170) (1.168311,0.666919) (1.172513,0.645144) (1.175353,0.622889) (1.176796,0.600200) (1.176808,0.577125) (1.175355,0.553712) (1.172403,0.530010) (1.167917,0.506069) (1.161864,0.481940) (1.154206,0.457677) (1.144907,0.433331) (1.133929,0.408957) (1.121232,0.384610) (1.106771,0.360346) (1.090499,0.336221) (1.072364,0.312295) (1.052307,0.288625) (1.030263,0.265273) (1.006154,0.242301) (0.979890,0.219772) (0.951365,0.197753) (0.920450,0.176310) (0.886984,0.155517) (0.850770,0.135447) (0.811553,0.116180) (0.769000,0.097804) (0.722661,0.080412) (0.671913,0.064111) (0.615843,0.049026) (0.553039,0.035305) (0.481111,0.023142) (0.395404,0.012806) (0.284114,0.004742) (0.069570,0.000070)};\draw[curve]plot coordinates{(0.069570,-0.000070) (0.284114,-0.004742) (0.395404,-0.012806) (0.481111,-0.023142) (0.553039,-0.035305) (0.615843,-0.049026) (0.671913,-0.064111) (0.722661,-0.080412) (0.769000,-0.097804) (0.811553,-0.116180) (0.850770,-0.135447) (0.886984,-0.155517) (0.920450,-0.176310) (0.951365,-0.197753) (0.979890,-0.219772) (1.006154,-0.242301) (1.030263,-0.265273) (1.052307,-0.288625) (1.072364,-0.312295) (1.090499,-0.336221) (1.106771,-0.360346) (1.121232,-0.384610) (1.133929,-0.408957) (1.144907,-0.433331) (1.154206,-0.457677) (1.161864,-0.481940) (1.167917,-0.506069) (1.172403,-0.530010) (1.175355,-0.553712) (1.176808,-0.577125) (1.176796,-0.600200) (1.175353,-0.622889) (1.172513,-0.645144) (1.168311,-0.666919) (1.162782,-0.688170) (1.155961,-0.708851) (1.147886,-0.728921) (1.138591,-0.748338) (1.128116,-0.767063) (1.116498,-0.785055) (1.103776,-0.802277) (1.089990,-0.818694) (1.075181,-0.834271) (1.059390,-0.848975) (1.042659,-0.862773) (1.025031,-0.875637) (1.006550,-0.887536) (0.987260,-0.898446) (0.967205,-0.908339) (0.946432,-0.917193) (0.924986,-0.924986) (0.902914,-0.931698) (0.880263,-0.937310) (0.857081,-0.941806) (0.833416,-0.945172) (0.809315,-0.947394) (0.784827,-0.948461) (0.760001,-0.948365) (0.734886,-0.947097) (0.709531,-0.944653) (0.683985,-0.941028) (0.658296,-0.936222) (0.632513,-0.930234) (0.606685,-0.923066) (0.580860,-0.914723) (0.555087,-0.905210) (0.529412,-0.894534) (0.503884,-0.882706) (0.478549,-0.869735) (0.453453,-0.855637) (0.428641,-0.840425) (0.404159,-0.824115) (0.380050,-0.806727) (0.356359,-0.788280) (0.333126,-0.768795) (0.310393,-0.748297) (0.288202,-0.726810) (0.266590,-0.704361) (0.245597,-0.680976) (0.225260,-0.656686) (0.205613,-0.631522) (0.186692,-0.605515) (0.168530,-0.578700) (0.151158,-0.551110) (0.134607,-0.522782) (0.118905,-0.493753) (0.104081,-0.464061) (0.090159,-0.433745) (0.077164,-0.402846) (0.065119,-0.371404) (0.054044,-0.339461) (0.043958,-0.307061) (0.034879,-0.274246) (0.026823,-0.241061) (0.019804,-0.207550) (0.013832,-0.173759) (0.008920,-0.139732) (0.005075,-0.105517) (0.002305,-0.071159) (0.000613,-0.036704) (0.000002,-0.002200)};\fill(0,0)circle(1.5pt);

\end{tikzpicture}
\caption{四极场的子午截面分为四个叶区；场线方程需连同角度范围使用。}\label{fig:es1-e9-q}
\end{figure}
\begin{solution}
在$\sin\theta\cos\theta\ne0$的区间内，切向比例为
\[\frac{\dd r}{r\dd\theta}
=\frac{3\cos^2\theta-1}{2\sin\theta\cos\theta}
=\cot\theta-\frac12\tan\theta.\]
在一个不穿越轴线和赤道的角区间内积分，
\[\int\frac{\dd r}r=\int\frac{\cos\theta}{\sin\theta}\dd\theta
-\frac12\int\frac{\sin\theta}{\cos\theta}\dd\theta
=\ln\sin\theta+\frac12\ln|\cos\theta|+C_0.\]
两次代换分别为$u=\sin\theta$和$v=\cos\theta$。因此$\ln r=\ln\sin\theta+\tfrac12\ln|\cos\theta|+C_0$，可写成
\[\boxed{r^2=C\sin^2\theta\cos\theta}.\]
北半球$C>0$，南半球$C<0$；也可统一写成$r=C_1\sin\theta\sqrt{|\cos\theta|}$，$C_1>0$。图~\ref{fig:es1-e9-s}用绝对值形式分别画两个半球，避免出现负的$r^2$。

球冠通量为
\[\begin{aligned}
\Psi&=2\pi\int_0^\theta E_r(r,u)r^2\sin u\dd u\\
&=-\frac{6\pi\ke A}{r^2}\int_{\cos\theta}^1(3v^2-1)\dd v\\
&=-\frac{6\pi\ke A}{r^2}[v^3-v]_{\cos\theta}^1
=-\frac{3A}{2\epsilon_0r^2}\sin^2\theta\cos\theta.
\end{aligned}\]
其中$v=\cos u$、$\dd v=-\sin u\dd u$；端点$u=0,\theta$分别变为$v=1,\cos\theta$。
在每个不穿过奇点的场管上$\Psi$不变，与上述隐式方程一致。轴线和赤道使$E_\theta=0$，是另行保留的径向特殊分支。

一条普通场线的$r$在$\cos^2\theta=1/3$处达到该叶区的最大值，此处$E_r=0$，但$E_\theta\ne0$，所以曲线沿球面方向转弯而不会终止。整个理想模型的原点为奇点；有限三电荷体系在源尺度附近的结构须用精确电势恢复。
\begin{figure}[H]\centering
\begin{tikzpicture}[>=Stealth,line cap=round,line join=round]
\draw[axis](-2.45,0)--(2.6,0)node[right]{$x$};\draw[axis](0,-2.05)--(0,2.2)node[above]{$z$};\draw[curve,domain=0.3:89.7,samples=100,smooth]plot ({-1.6*(sin(\x))^2*sqrt(abs(cos(\x)))},{1.6*sin(\x)*sqrt(abs(cos(\x)))*cos(\x)});\draw[curve,domain=90.3:179.7,samples=100,smooth]plot ({-1.6*(sin(\x))^2*sqrt(abs(cos(\x)))},{1.6*sin(\x)*sqrt(abs(cos(\x)))*cos(\x)});\draw[->](-0.49721,0.68435)--(-0.37224,0.64474);\draw[->](-0.75946,0.20350)--(-0.83481,0.32045);\draw[->](-0.75946,-0.20350)--(-0.83481,-0.32045);\draw[->](-0.49721,-0.68435)--(-0.37224,-0.64474);\draw[curve,domain=0.3:89.7,samples=100,smooth]plot ({1.6*(sin(\x))^2*sqrt(abs(cos(\x)))},{1.6*sin(\x)*sqrt(abs(cos(\x)))*cos(\x)});\draw[curve,domain=90.3:179.7,samples=100,smooth]plot ({1.6*(sin(\x))^2*sqrt(abs(cos(\x)))},{1.6*sin(\x)*sqrt(abs(cos(\x)))*cos(\x)});\draw[->](0.49721,0.68435)--(0.37224,0.64474);\draw[->](0.75946,0.20350)--(0.83481,0.32045);\draw[->](0.75946,-0.20350)--(0.83481,-0.32045);\draw[->](0.49721,-0.68435)--(0.37224,-0.64474);\draw[curve,domain=0.3:89.7,samples=100,smooth]plot ({-2.5*(sin(\x))^2*sqrt(abs(cos(\x)))},{2.5*sin(\x)*sqrt(abs(cos(\x)))*cos(\x)});\draw[curve,domain=90.3:179.7,samples=100,smooth]plot ({-2.5*(sin(\x))^2*sqrt(abs(cos(\x)))},{2.5*sin(\x)*sqrt(abs(cos(\x)))*cos(\x)});\draw[->](-0.77688,1.06929)--(-0.58163,1.00741);\draw[->](-1.18666,0.31796)--(-1.30439,0.50071);\draw[->](-1.18666,-0.31796)--(-1.30439,-0.50071);\draw[->](-0.77688,-1.06929)--(-0.58163,-1.00741);\draw[curve,domain=0.3:89.7,samples=100,smooth]plot ({2.5*(sin(\x))^2*sqrt(abs(cos(\x)))},{2.5*sin(\x)*sqrt(abs(cos(\x)))*cos(\x)});\draw[curve,domain=90.3:179.7,samples=100,smooth]plot ({2.5*(sin(\x))^2*sqrt(abs(cos(\x)))},{2.5*sin(\x)*sqrt(abs(cos(\x)))*cos(\x)});\draw[->](0.77688,1.06929)--(0.58163,1.00741);\draw[->](1.18666,0.31796)--(1.30439,0.50071);\draw[->](1.18666,-0.31796)--(1.30439,-0.50071);\draw[->](0.77688,-1.06929)--(0.58163,-1.00741);\draw[curve,domain=0.3:89.7,samples=100,smooth]plot ({-3.4*(sin(\x))^2*sqrt(abs(cos(\x)))},{3.4*sin(\x)*sqrt(abs(cos(\x)))*cos(\x)});\draw[curve,domain=90.3:179.7,samples=100,smooth]plot ({-3.4*(sin(\x))^2*sqrt(abs(cos(\x)))},{3.4*sin(\x)*sqrt(abs(cos(\x)))*cos(\x)});\draw[->](-1.05656,1.45423)--(-0.79101,1.37008);\draw[->](-1.61386,0.43243)--(-1.77397,0.68096);\draw[->](-1.61386,-0.43243)--(-1.77397,-0.68096);\draw[->](-1.05656,-1.45423)--(-0.79101,-1.37008);\draw[curve,domain=0.3:89.7,samples=100,smooth]plot ({3.4*(sin(\x))^2*sqrt(abs(cos(\x)))},{3.4*sin(\x)*sqrt(abs(cos(\x)))*cos(\x)});\draw[curve,domain=90.3:179.7,samples=100,smooth]plot ({3.4*(sin(\x))^2*sqrt(abs(cos(\x)))},{3.4*sin(\x)*sqrt(abs(cos(\x)))*cos(\x)});\draw[->](1.05656,1.45423)--(0.79101,1.37008);\draw[->](1.61386,0.43243)--(1.77397,0.68096);\draw[->](1.61386,-0.43243)--(1.77397,-0.68096);\draw[->](1.05656,-1.45423)--(0.79101,-1.37008);\fill(0,0)circle(1.4pt);\node[right,align=left]at(2.95,.85){北半球$C>0$\\[5pt]南半球$C<0$\\[5pt]$r^2=C\sin^2\theta\cos\theta$};
\end{tikzpicture}
\caption{四极场线在北半球沿极角减小方向、南半球沿极角增大方向行进。}\label{fig:es1-e9-s}
\end{figure}


图~\ref{fig:es1-e9-s}中，北半球$E_\theta<0$，南半球$E_\theta>0$，箭头因而都由赤道附近指向轴向源区。轴上场指向负电荷源区，赤道上的场沿径向向外，这也确定了两条特殊分支的方向。
\end{solution}

\begin{example}[零势面是否也是零场面]\label{ex:es1-e10}
对理想偶极和例~\ref{ex:es1-e8}的理想四极子，分别求零势面的方向，并判断这些方向上电场是否为零。 见图~\ref{fig:es1-e10-q}。
\end{example}

\begin{figure}[H]\centering
\begin{tikzpicture}[>=Stealth,line cap=round,line join=round]

\draw[axis](-1.7,0)--(1.9,0)node[right]{$x$};\draw[axis](0,-1.7)--(0,1.9)node[above]{$z$};
\draw[dashed](-1.8,-1.27279)--(1.8,1.27279) (-1.8,1.27279)--(1.8,-1.27279);\draw[curve](-1.6,0)--(1.6,0);\node[right]at(2.5,.6){$\varphi=0$的候选方向};

\end{tikzpicture}
\caption{偶极的赤道面与四极的两张零势锥面在截面上呈现不同几何。}\label{fig:es1-e10-q}
\end{figure}
\begin{solution}
偶极势正比$\cos\theta$，故零势面为$\theta=\pi/2$。此处$E_r=0$，但$E_\theta=\ke p/r^3$，电场沿负偶极轴，不为零。

四极势正比$3\cos^2\theta-1$，故两张零势锥满足$\cos\theta=\pm1/\sqrt3$。在锥上$E_r=0$，而
\[|E_\theta|=\frac{6\ke A}{r^4}\frac1{\sqrt3}\sqrt{\frac23}
=\frac{2\sqrt2\ke A}{r^4}>0.\]
图~\ref{fig:es1-e10-s}在北半球锥面上一点画出沿$-\vct e_\theta$的场。由于电势在整个锥面为零，场应当垂直这个等势面；对固定极角的锥面，法向恰是角向方向。

把电势加任意常数还会改变哪一张面叫作零势面，却不会改变电场。零势面取决于势零点的选择，零场点由梯度是否为零决定。真正的零场点要求全部独立分量同时为零；在理想偶极和本四极模型中，除无穷远外不存在这样的普通点。有限源体系则可能有额外的零场点，不能仅由远区式排除。
\begin{figure}[H]\centering
\begin{tikzpicture}[>=Stealth,line cap=round,line join=round]

\draw[axis](-.4,0)--(2.4,0)node[right]{$x$};\draw[axis](0,-.3)--(0,2.4)node[above]{$z$};
\draw[curve](0,0)--(2.2,1.55563);\fill(1.2,.848528)circle(1.5pt)node[below right]{$P$};
\draw[axis](1.2,.848528)--(.65,1.62635)node[above right]{$\vct E\parallel-\vct e_\theta$};\draw[guide](0,0)--(1.2,.848528);

\end{tikzpicture}
\caption{四极零势锥面上的场沿角向，既不沿锥面母线，也不等于零。}\label{fig:es1-e10-s}
\end{figure}


四极锥的半顶角为$\theta_0=\arccos(1/\sqrt3)$，截面母线满足$z=\pm x/\sqrt2$。在北半球、$x>0$的母线上，$E_\theta<0$，所以场指向$-\vct e_\theta$，即向上并靠近轴。母线斜率$1/\sqrt2$与场箭头斜率$-\sqrt2$相乘为$-1$。用这个点积关系可以直接核对图中箭头是否垂直等势锥面。
\end{solution}

\subsection{电场能、组装过程与表面项}\label{sec:es-field-energy}
先考虑有界、光滑的真空电荷分布。将分布从零按比例建立：某时刻源为$t\rho$，电势为$t\varphi$；增加$\rho\dd t$所需的功为
\[\dd W=\int(t\varphi)\rho\dd V\,\dd t.
\]
从$t=0$积到$1$，得到
\[W=\left(\int_0^1t\dd t\right)\int\rho\varphi\dd V
=\left[\frac{t^2}2\right]_0^1\int\rho\varphi\dd V
=\frac12\int\rho\varphi\dd V.
\]
系数$1/2$来自势随电量同步增加。对两个预先给定的点源，相互作用能仍为$\ke q_1q_2/r_{12}$，因为求和式中的两项相同：$\tfrac12(q_1\varphi_2+q_2\varphi_1)=\ke q_1q_2/r_{12}$。这里没有计入点源自身的奇异能量。

将$\rho=\epsilon_0\nabla\cdot\vct E$代入，并使用
\[\nabla\cdot(\varphi\vct E)
=\nabla\varphi\cdot\vct E+\varphi\nabla\cdot\vct E
=-E^2+\frac{\rho\varphi}{\epsilon_0},
\]
定义有限区域内的源能积分$W_V=\tfrac12\int_V\rho\varphi\dd V$，再对上式在$V$内积分，便有
\[W_V=\frac{\epsilon_0}{2}\int_VE^2\dd V
+\frac{\epsilon_0}{2}\oint_{\partial V}\varphi\vct E\cdot\vct n\dd S.
\]
对有限总电量的孤立体系，远处$\varphi=O(r^{-1})$、$E=O(r^{-2})$，无穷远表面项按$r^{-1}$趋零，于是得到正的场能密度$w=\epsilon_0E^2/2$。若只取有限区域、外场不衰减或边界接有电源，表面项一般不能删去。点电荷附近还须先挖去小球，分开处理发散的自能；它与有限的点源互能是不同的量。

导体与介质的诱导分布随外源改变时，组装中的半因子也会出现。第三讲的电像能、第四讲的诱导偶极能都由同一个线性建立过程得到。受力时还须说明固定的是电量还是电压：电源与体系交换的功不能并入一个不加区分的场能导数。



\section{静电关系在完整场方程中的位置}\label{sec:es1-maxwell}
静电条件是完整场方程的一种时间不变情形。将自由源与束缚源分开后，可以看清静电、稳恒电流和静磁场各自保留的方程及额外条件。

在介质中分离自由与束缚源，定义$\vct D=\epsilon_0\vct E+\vct P$、$\vct B=\mu_0(\vct H+\vct M)$。宏观场方程的微分形式为
\[\nabla\cdot\vct D=\rho_f,\qquad \nabla\cdot\vct B=0,\]
\[\nabla\times\vct E=-\partial_t\vct B,\qquad
\nabla\times\vct H=\vct j_f+\partial_t\vct D.\]
对固定体积$V$和固定有向曲面$S$积分，分别使用散度定理和Stokes定理，得
\[\oint_{\partial V}\vct D\cdot\vct n\dd S=\int_V\rho_f\dd V,
\qquad \oint_{\partial V}\vct B\cdot\vct n\dd S=0,\]
\[\oint_{\partial S}\vct E\cdot\dd\vct l
=-\frac{\dd}{\dd t}\int_S\vct B\cdot\vct n\dd S,\]
\[\oint_{\partial S}\vct H\cdot\dd\vct l
=\int_S\vct j_f\cdot\vct n\dd S
+\frac{\dd}{\dd t}\int_S\vct D\cdot\vct n\dd S.\]
闭面的法向取向外；开曲面的法向与边界绕行按右手规则配对。曲面固定才可将时间导数移到通量积分外。若回路运动，还需计入运动电动势，不能将上述固定面公式直接换上移动边界。

静电情形的$\partial_t\vct B=0$使电场环量为零。静磁情形则有$\nabla\times\vct H=\vct j_f$。真空中$\vct D=\epsilon_0\vct E$、$\vct H=\vct B/\mu_0$，才可把电荷写成$\epsilon_0$乘电场通量、把电流写成磁感应强度环量除以$\mu_0$。第四讲进一步说明介质中这些量与源的对应关系。


固定回路的环量、场的局部旋度和电势的存在性各有其条件。静电区域中若$\partial_t\vct B=0$，则小回路环量为零，可在单连通域定义$\vct E=-\nabla\varphi$。若区域包围奇异轴或穿孔，局部无旋还不足以保证全域单值势，第一讲有关角度势的习题正对应这种情形。只在一条曲线上知道势差，也不足以重建附近的横向场；场需要完整空间梯度。

位移电流与电荷连续性可直接从方程联系起来。对$\nabla\times\vct H=\vct j_f+\partial_t\vct D$取散度，左边恒为零，故
\[0=\nabla\cdot\vct j_f+\partial_t(\nabla\cdot\vct D)
=\nabla\cdot\vct j_f+\partial_t\rho_f.
\]
在稳恒导电问题中$\partial_t\rho_f=0$，电流无散度；这并不等于电场无散度。导电率或介电常数在空间变化时，维持稳恒电流仍可能需要体电荷与界面电荷。第三讲讨论的绝缘孔要满足电流不穿壁，第四讲的介质界面则按自由电荷规定$D$的法向跳跃，二者应从各自守恒的通量建立边界。

对线性、各向同性的静电介质，$\vct D=\epsilon(\vct r)\vct E$，势方程为
\[\nabla\cdot[\epsilon(\vct r)\nabla\varphi]=-\rho_f.
\]
只有在一个子区域中$\epsilon$为常数且$\rho_f=0$，才可简化为$\nabla^2\varphi=0$。介质分界上仍须另接势连续和法向位移条件。各子域的Laplace解只有同时满足界面条件，才构成整个区域的解。



\providecommand{\ExpansionDiagram}[3]{\begin{figure}[H]\centering\begin{tikzpicture}#3\end{tikzpicture}\caption{#2}\label{#1}\end{figure}}
\section{由场能到导体表面的电压力}
\subsection{面电荷所受的场不是单侧总场}
真空中一层面电荷密度$\sigma$的薄层，两侧法向场之差为$\sigma/\varepsilon_0$。把局部场分成该层自身的奇对称场和其余电荷的平滑场，薄层自身场在两侧分别为$\pm\sigma/(2\varepsilon_0)$，因而其余电荷在薄层处的场是两侧总场的平均值：
\[
 \vct E_{\rm other}=\frac{\vct E_++\vct E_-}{2}.
\]
单独取外侧场乘$\sigma$，会把该层对自身的作用重复计入。对静电平衡导体，内侧场为零、外侧场为$E_n\vct n$，于是单位面积的电力为
\[
 \vct f_s=\sigma\frac{E_n}{2}\vct n
 =\frac{\varepsilon_0 E_n^2}{2}\vct n
 =\frac{\sigma^2}{2\varepsilon_0}\vct n.
\]
无论局部电荷正负，电压力均沿导体外法线。若表面两侧都是真空场，应分别计算两侧的应力，取它们对同一薄层的合力，而非只取某一侧。
\ExpansionDiagram{fig:dep-es1-pillbox}{导体表面薄层；外侧场与其余电荷场相差一倍}{
\fill[gray!15](-3,-1.2)rectangle(3,0);\draw[thick](-3,0)--(3,0);\node at(-1.5,-.8){导体：$E_-=0$};\draw[dashed](-.5,-.3)rectangle(.5,.6);\draw[->,thick](0,.2)--(0,1.8)node[above]{$E_n$};\draw[->](1.7,.2)--(1.7,1)node[above]{$E_n/2$};\node[right]at(1.9,.5){其余源};}
图\ref{fig:dep-es1-pillbox}中虚线薄盒包住一小块表面。$\varepsilon_0E_n^2/2$也正是真空电场能量密度，但在这里它具有法向压力的力学含义，量纲为$\mathrm{J/m^3=Pa}$。

\subsection{应力张量怎样给任意封闭面的总力}
若闭合面穿过真空且包住全部受力电荷，可定义静电应力
\[
 T_{ij}=\varepsilon_0(E_iE_j-\tfrac12E^2\delta_{ij}),
 \qquad \vct F=\oint \mathsf T\vct n\,\dd S.
\]
其局部依据可逐分量求出，重复指标表示求和。静电无旋给$\partial_jE_i=\partial_iE_j$，故
\[E_j\partial_jE_i=E_j\partial_iE_j
=\frac12\partial_i(E_jE_j)=\frac12\partial_iE^2.\]
对$T_{ij}$使用乘积法则，再代入$\partial_jE_j=\rho/\varepsilon_0$，
\[\begin{aligned}
\partial_jT_{ij}
&=\varepsilon_0\left[(\partial_jE_i)E_j
+E_i\partial_jE_j-\frac12\partial_iE^2\right]\\
&=\varepsilon_0E_i\partial_jE_j=\rho E_i.
\end{aligned}\]
前后两个场强平方的导数相消，余下的是电荷的力密度。再用散度定理，即得闭面积分公式。积分面可选在场较简单的真空区域。

若积分面贴近导体，$\vct E=E_n\vct n$，所以$\mathsf T\vct n=\varepsilon_0E_n^2\vct n/2$，恢复局部压力公式。这里使用真空电磁应力；在可压缩介质中，还要交代材料应力、极化和电致伸缩的模型，不能无条件把$\varepsilon_0$换成$\varepsilon$。

\begin{example}[带电导体球的局部压力与整体受力]
半径$a$的导体球带总电量$Q$，置于沿$+z$的均匀外场$E_0$。求表面电场和压力分布、整个球的合力，以及外场单独造成的北半球轴向电力。忽略支承结构对场的影响。
\ExpansionDiagram{fig:dep-es1-sphere-q}{带电球与轴向外场；北半球取$0\leq\theta\leq\pi/2$}{
\draw[thick](0,0)circle(1.3);\draw[->](0,-1.8)--(0,2.6)node[above]{$z,E_0$};\draw[->](0,0)--(.92,.92)node[midway,below right]{$a$};\draw(0,.45)arc(90:45:.45);\node at(.27,.65){$\theta$};\draw[dashed](-1.3,0)--(1.3,0);\node[left]at(-1.4,.6){北半球};\node at(.5,-.6){$Q$};}
\end{example}
\begin{solution}
在球外把单极势与匀强场中的导体球解叠加：
\[
 \phi=\frac{Q}{4\pi\varepsilon_0r}
 -E_0\left(r-\frac{a^3}{r^2}\right)\cos\theta.
\]
在$r=a$角向项为零，球面等势；远处恢复原外场，单极项保证总电量为$Q$，因此满足题设。先保持$\theta$不变作径向微分，
\[-\partial_r\phi=\frac{Q}{4\pi\varepsilon_0r^2}
+E_0\left(1+\frac{2a^3}{r^3}\right)\cos\theta.\]
再令$r=a$，得到
\[
 E_n=\frac{Q}{4\pi\varepsilon_0a^2}+3E_0\cos\theta
 \equiv E_Q+3E_0\cos\theta.
\]
图\ref{fig:dep-es1-sphere-q}中的北半球外法线有正$z$分量$\cos\theta$，故压力的轴向贡献为$\varepsilon_0(E_Q+3E_0\cos\theta)^2\cos\theta/2$。
\ExpansionDiagram{fig:dep-es1-sphere-a}{局部压力沿径向；整体轴向力来自交叉项的不对称}{
\draw[thick](0,0)circle(1.2);\foreach\a in{30,60,90,120,150,210,240,270,300,330}{\draw[->](\a:1.2)--(\a:{1.2+.015*(4+3*sin(\a))^2});}\draw[->](2,-.6)--(2,1.2)node[above]{$Q E_0$};\node[below]at(0,-1.9){图示$E_Q=4E_0$，箭长按压力比例};}
对整个球积分，令$u=\cos\theta$，方位角积分为$2\pi$：
\begin{align*}
 F_z&=\pi\varepsilon_0a^2\int_{-1}^1(E_Q+3E_0u)^2u\dd u\\
 &=\pi\varepsilon_0a^2\left(6E_QE_0\int_{-1}^1u^2\dd u\right)
 =6\pi\varepsilon_0a^2E_QE_0\left[\frac{u^3}{3}\right]_{-1}^1
 =4\pi\varepsilon_0a^2E_QE_0=QE_0.
\end{align*}
纯自场的$E_Q^2u$项和纯外场诱导的$9E_0^2u^3$项为奇函数，积分为零；合力来自交叉项。图\ref{fig:dep-es1-sphere-a}仅示意$Q>0$时的压力不对称，$Q<0$时合力反向，局部压力仍向外。

若只取北半球且令$Q=0$，仍有非零轴向电力：
\[
 F_{N,z}=9\pi\varepsilon_0a^2E_0^2\int_0^1u^3\dd u
 =9\pi\varepsilon_0a^2E_0^2\left[\frac{u^4}{4}\right]_0^1
 =\frac94\pi\varepsilon_0a^2E_0^2.
\]
南半球受到大小相同、方向相反的电力，整个球的合力为零，两个半球仍需由内部机械应力维持连接。若要求变形量，还应给出材料弹性或表面张力，把电压力作为载荷求解。
\end{solution}

\section{习题与解答}

\begin{exercise}[球壳上的势与场]\label{ex:es1-h1}
半径$a$的均匀带电薄球壳，总电荷$Q>0$。取无穷远零势，求内外电场和电势，并说明球面上的哪一个场值能用于求壳上一小块电荷的受力。 见图~\ref{fig:es1-h1-q}。
\end{exercise}

\begin{figure}[H]\centering
\begin{tikzpicture}[>=Stealth,line cap=round,line join=round]
\draw[curve](0,0)circle(1.5);\fill(0,0)circle(1.4pt)node[below left]{$O$};\node at(-.7,.8){$Q$};\draw[axis](0,0)--(1.5,0)node[midway,above]{$a$};
\end{tikzpicture}
\caption{球壳的总电荷与几何半径均已给定。}\label{fig:es1-h1-q}
\end{figure}
\begin{solution}
球对称使场仅有径向分量。半径$r<a$的Gauss球不含电荷，故$E=0$；$r>a$时$E=\ke Q/r^2$。从无穷远沿径向积分，
\[\varphi(r)=-\int_\infty^r\frac{\ke Q}{s^2}\dd s
=-\ke Q\left[-\frac1s\right]_\infty^r=\frac{\ke Q}{r}\quad(r>a).\]
内区$\varphi(r)-\varphi(a)=-\int_a^r0\dd s=0$。使势连续，
\[\varphi=\begin{cases}\ke Q/a,&r\le a,\\\ke Q/r,&r\ge a.\end{cases}\]
图~\ref{fig:es1-h1-s}中势连续，导数在球面跳变。求壳上一小块的受力时，须去掉这块电荷自身的局部场；对理想薄层，其他电荷的场为两侧极限的平均$E_{\rm other}=(E_{\rm out}+E_{\rm in})/2=\ke Q/(2a^2)$，方向向外。

因此不能把$E_{\rm out}$全乘到同一带电面元上，那会多计自场的一半。对远处试探电荷则使用它所在一侧的实际场，不取两侧平均。
\begin{figure}[H]\centering
\begin{tikzpicture}[>=Stealth,line cap=round,line join=round]
\begin{scope}[x=1.4cm,y=1.3cm]\draw[axis](0,0)--(3.18,0)node[right]{$r/a$};\draw[axis](0,0)--(0,1.4)node[above]{$\varphi/(\ke Q/a)$};\draw[curve,domain=1:3,samples=101]plot(\x,{1/\x});\draw[curve](0,1)--(1,1);\draw[guide](1,0)--(1,1);\node[below]at(1,0){$1$};\end{scope}

\end{tikzpicture}
\caption{壳内电势为常数，壳外按$1/r$下降；电场跃变由自由面电荷造成。}\label{fig:es1-h1-s}
\end{figure}


由$\sigma=\epsilon_0E_{\rm out}$，薄壳的面力密度为
\[\vct f_s=\sigma E_{\rm other}\vct e_r
=\frac12\epsilon_0E_{\rm out}^2\vct e_r,\]
沿外法向。这一压力使用扣除面元自场后的平均场；壳外试探电荷则使用所在处的实际电场。
\end{solution}

\begin{exercise}[从环轴到平面]\label{ex:es1-h2}
半径$a$的均匀细圆环总电荷$Q$。求轴线上$z$处的电势与场，再用环元叠加求半径$R$、面密度$\sigma$的圆盘轴向场，讨论无限平面极限。 见图~\ref{fig:es1-h2-q}。
\end{exercise}

\begin{figure}[H]\centering
\begin{tikzpicture}[>=Stealth,line cap=round,line join=round]
\draw[curve](0,0)ellipse(1.6 and .45);\draw[axis](0,-.6)--(0,2.5)node[above]{$z$};\fill(0,1.8)circle(1.5pt)node[right]{$P$};\draw[guide](0,1.8)--(1.6,0);\node[above]at(.8,0){$a$};
\end{tikzpicture}
\caption{同一环上的电荷到轴上点等距，横向场成对抵消。}\label{fig:es1-h2-q}
\end{figure}
\begin{solution}
环上各元到观察点距离相同，所以$\varphi=\ke\int\dd q/\sqrt{a^2+z^2}=\ke Q/\sqrt{a^2+z^2}$。链式求导给出
\[E_z=-\ke Q\left(-\frac12\right)(a^2+z^2)^{-3/2}(2z)
=\frac{\ke Qz}{(a^2+z^2)^{3/2}}.\]
对圆盘用半径$s$、电量$\dd Q=2\pi\sigma s\dd s$的环元叠加。令$w=s^2+z^2$，则$\dd w=2s\dd s$，端点为$z^2,R^2+z^2$。$z>0$时
\[\begin{aligned}
E_z&=2\pi\ke\sigma\int_0^R\frac{zs\dd s}{(s^2+z^2)^{3/2}}\\
&=\pi\ke\sigma z\int_{z^2}^{R^2+z^2}w^{-3/2}\dd w\\
&=\pi\ke\sigma z[-2w^{-1/2}]_{z^2}^{R^2+z^2}
=\frac\sigma{2\epsilon_0}\left(1-\frac z{\sqrt{R^2+z^2}}\right).
\end{aligned}\]
图~\ref{fig:es1-h2-s}的环元同时决定积分权重与距离，漏掉$2\pi s$会使量纲错误。下方场按对称性反号。固定非零$z$令$R\to\infty$，得到两侧大小均为$\sigma/(2\epsilon_0)$的平面场。

无限平面的场不衰减，不能再把无穷远作为有限零势。圆盘势在有限$R$时可取无穷远零势，但其$R\to\infty$极限含一个发散常数；势差和场的极限仍可存在。
\begin{figure}[H]\centering
\begin{tikzpicture}[>=Stealth,line cap=round,line join=round]
\draw[curve](0,0)ellipse(1.7 and .5);\draw[thick](0,0)ellipse(1.05 and .31);\draw[guide](0,0)--(1.05,0)node[midway,above]{$s$};\draw[axis](0,0)--(0,2.3)node[above]{$z$};\node[right]at(2.2,.7){$\dd Q=2\pi\sigma s\dd s$};
\end{tikzpicture}
\caption{圆盘分解为同轴带电环，积分上下限为$0$与$R$。}\label{fig:es1-h2-s}
\end{figure}


圆盘轴上势由同一环元积分得到：
\[\varphi(z)=2\pi\ke\sigma\int_0^R\frac{s\dd s}{\sqrt{s^2+z^2}}
=\pi\ke\sigma[2\sqrt w]_{z^2}^{R^2+z^2}
=\frac{\sigma}{2\epsilon_0}(\sqrt{R^2+z^2}-|z|).\]
对$z\ne0$，$\dd\sqrt{R^2+z^2}/\dd z=z/\sqrt{R^2+z^2}$，$\dd|z|/\dd z=\operatorname{sgn}z$，故取负导数得到两侧方向相反的场。$z=0$处势连续但导数跃变，$E_z(0^+)-E_z(0^-)=\sigma/\epsilon_0$。当$z\gg R$，展开根号得到$\varphi\simeq\sigma\pi R^2/(4\pi\epsilon_0z)$，正好是总电荷的点源势。
\end{solution}

\begin{exercise}[半球面的固体角]\label{ex:es1-h3}
点电荷$Q$在球心。求球冠$0\le\theta\le\alpha$的外通量，再令点电荷沿轴移到球面上一点的内侧极限，说明不能继续使用均匀场乘球冠面积。 见图~\ref{fig:es1-h3-q}。
\end{exercise}

\begin{figure}[H]\centering
\begin{tikzpicture}[>=Stealth,line cap=round,line join=round]
\draw[curve](0,0)circle(1.5);\fill(0,0)circle(1.4pt)node[below left]{$O$};\draw[very thick](1.2,.9)arc(36.87:143.13:1.5);\draw[guide](-1.2,.9)--(1.2,.9);\draw[axis](0,0)--(0,2)node[above]{$z$};\draw[guide](0,0)--(1.2,.9);\draw(0,.5)arc(90:36.87:.5);\node at(.25,.62){$\alpha$};\node[below left]at(0,0){$Q$};

\end{tikzpicture}
\caption{第一种构形中，所有球面点到电荷等距。}\label{fig:es1-h3-q}
\end{figure}
\begin{solution}
球心电荷的径向场在球面恒定，球冠面积为$2\pi R^2(1-\cos\alpha)$，故
\[\Phi=\frac{Q}{2\epsilon_0}(1-\cos\alpha).\]
它也等于$Q\Omega/(4\pi\epsilon_0)$，固体角$\Omega=2\pi(1-\cos\alpha)$。

图~\ref{fig:es1-h3-s}把偏离球心后的观察射线画出。电荷接近球面时，各面元到电荷的距离和入射角都变化，不能仍用球心构形的面积公式。应从电荷位置观察该曲面的有向固体角，或逐面积积分。若对整个闭球面求通量，只要电荷仍在内部就保持$Q/\epsilon_0$；穿过边界后变为零。边界上的理想点电荷本身是奇点，需先说明取哪一侧极限。
\begin{figure}[H]\centering
\begin{tikzpicture}[>=Stealth,line cap=round,line join=round]
\draw[curve](0,0)circle(1.5);\fill(0,0)circle(1.4pt)node[below left]{$O$};\fill(0,1.2)circle(1.7pt)node[left]{$Q$};\draw[guide](0,1.2)--(1.4,.5385) (0,1.2)--(-1.2,-.9);\draw[axis](1.4,.5385)--(1.85,.7116)node[right]{$\vct n$};
\end{tikzpicture}
\caption{偏心源点的射线与球面法向不再平行，通量须包含投影因子。}\label{fig:es1-h3-s}
\end{figure}


偏心电荷位于轴上$z=d<R$时，记$L(\theta)=\sqrt{R^2+d^2-2Rd\cos\theta}$。面元法向投影为$(R-d\cos\theta)/L(\theta)$，所以
\[\Phi=2\pi\ke Q\int_0^\alpha\frac{R^2(R-d\cos\theta)\sin\theta}{L(\theta)^3}\dd\theta.\]
为求原函数，对边缘射线的方向余弦求导：
\[\frac{\dd}{\dd\theta}\frac{R\cos\theta-d}{L(\theta)}
=-\frac{R\sin\theta}{L}-\frac{(R\cos\theta-d)Rd\sin\theta}{L^3}
=-\frac{R^2(R-d\cos\theta)\sin\theta}{L^3}.\]
故
\[\Phi=-2\pi\ke Q\left[\frac{R\cos\theta-d}{L(\theta)}\right]_0^\alpha
=\frac{Q}{2\epsilon_0}\left[1-\frac{R\cos\alpha-d}{L(\alpha)}\right].\]
这里$L(0)=R-d>0$。相应有向固体角为$\Omega=2\pi[1-(R\cos\alpha-d)/L(\alpha)]$。对固定$0<\alpha<\pi$取$d\uparrow R$，分式趋于$-\sin(\alpha/2)$，故$\Omega\to2\pi[1+\sin(\alpha/2)]$。尖端附近的有限通量在极限中集中，不能把源直接放到边界上再沿用普通点的积分公式。
\end{solution}

\begin{exercise}[不等量异号点电荷]\label{ex:es1-h4}
$+4Q$在$x=0$，$-Q$在$x=d>0$。求轴上的零场点与零势点，并说明两类点为何不同。 见图~\ref{fig:es1-h4-q}。
\end{exercise}

\begin{figure}[H]\centering
\begin{tikzpicture}[>=Stealth,line cap=round,line join=round]
\draw[axis](-2.7,0)--(2.9,0)node[right]{$x$};\fill(-1.2,0)circle(1.6pt)node[below]{$+4Q$};\fill(1.2,0)circle(1.6pt)node[below]{$-Q$};
\end{tikzpicture}
\caption{先将轴分成两源左侧、两源之间和两源右侧三个区间。}\label{fig:es1-h4-q}
\end{figure}
\begin{solution}
两源之间的场均向右，不可能为零。左侧若$|x|=u>0$，场抵消要求$4/u^2=1/(u+d)^2$，无正解。右侧$x>d$时方向相反，$4/x^2=1/(x-d)^2$给出$x=2d$，所以零场点只有$2d$。

零势要求$4/|x|=1/|x-d|$，分别在两源之间和右侧得到$x=4d/5$、$x=4d/3$。图~\ref{fig:es1-h4-s}沿同一轴标出三个位置。零势点可以有非零斜率，即非零场；零场点则是势的驻点，势值未必为零。

本题总电荷为$3Q$，远处势仍为正的单极首项。两个异号源之间虽有一个零势点，却没有电场反向的零点。按区间先判方向，可排除平方运算产生的伪解。
\begin{figure}[H]\centering
\begin{tikzpicture}[>=Stealth,line cap=round,line join=round]
\draw[axis](-.2,0)--(6.4,0)node[right]{$x/d$};\foreach\x/\lab in{0/0,3/1,2.4/{4/5},4/{4/3},6/2}{\draw(\x,-.08)--(\x,.08);\node[below]at(\x,0){$\lab$};}\node[above]at(2.4,.2){零势};\node[above]at(4,.2){零势};\node[above]at(6,.2){零场};
\end{tikzpicture}
\caption{势零点与场零点分别来自距离的一次比例和平方反比。}\label{fig:es1-h4-s}
\end{figure}


从势的一维曲线还能区分两类根。在$x>d$区，$\varphi=\ke Q[4/x-1/(x-d)]$，连续求导得
\[\begin{aligned}
\varphi'&=\ke Q[-4x^{-2}+(x-d)^{-2}],\\
\varphi''&=\ke Q[8x^{-3}-2(x-d)^{-3}].
\end{aligned}\]
在$x=2d$处，$\varphi'=0$、$\varphi=\ke Q/d>0$，而$\varphi''=\ke Q(d^{-3}-2d^{-3})=-\ke Q/d^3<0$。正试探电荷的势能沿轴取极大值。完整三维稳定性还须同时考察横向曲率。
\end{solution}

\begin{exercise}[偶极轴上与赤道上的远场]\label{ex:es1-h5}
有限电荷对$\pm q$相距$\ell$。求轴上与赤道上距中心$r$处的精确场，并比较偶极近似的首个修正项，取$r>\ell/2$。 见图~\ref{fig:es1-h5-q}。
\end{exercise}

\begin{figure}[H]\centering
\begin{tikzpicture}[>=Stealth,line cap=round,line join=round]
\draw[curve](0,0)circle(1.5);\fill(0,0)circle(1.4pt)node[above right]{$O$};\fill(0,.35)circle(1.6pt)node[left]{$+q$};\fill(0,-.35)circle(1.6pt)node[left]{$-q$};\fill(0,1.5)circle(1.5pt)node[right]{$P_a$};\fill(1.5,0)circle(1.5pt)node[above]{$P_e$};

\end{tikzpicture}
\caption{两个观察方向使用相同中心距离，但到两电荷的几何不同。}\label{fig:es1-h5-q}
\end{figure}
\begin{solution}
在正轴上两场反向，
\[E_a=\ke q[(r-\ell/2)^{-2}-(r+\ell/2)^{-2}]
=\frac{2\ke pr}{(r^2-\ell^2/4)^2}
=\frac{2\ke p}{r^3}\left[1+\frac{\ell^2}{2r^2}+O(\ell^4/r^4)\right].\]
赤道上横向分量抵消，轴向分量同向指向负电荷，大小
\[E_e=\frac{\ke p}{(r^2+\ell^2/4)^{3/2}}
=\frac{\ke p}{r^3}\left[1-\frac{3\ell^2}{8r^2}+O(\ell^4/r^4)\right].\]
图~\ref{fig:es1-h5-s}分别标出加减关系。远区轴向场与赤道场大小之比趋于$2$，但有限距离下不严格等于$2$。轴上修正增大场，赤道修正减小场；同一误差阶并不意味着各方向误差同号。
\begin{figure}[H]\centering
\begin{tikzpicture}[>=Stealth,line cap=round,line join=round]
\draw[axis](0,-1.2)--(0,2.1)node[above]{$z$};\draw[axis](0,0)--(2.6,0)node[right]{$x$};\draw[axis](0,1.3)--(0,1.9)node[right]{$\vct E_a$};\draw[axis](1.5,0)--(1.5,-.8)node[right]{$\vct E_e$};\fill(0,.3)circle(1.5pt);\fill(0,-.3)circle(1.5pt);
\end{tikzpicture}
\caption{两处合场都沿偶极轴，正轴上向上，赤道上向下。}\label{fig:es1-h5-s}
\end{figure}


以上修正系数可统一用$u=\ell^2/(4r^2)$求出。由二项式展开，
\[(1-u)^{-2}=1+2u+O(u^2),\qquad
(1+u)^{-3/2}=1-\frac32u+O(u^2),\]
分别乘轴向和赤道场的首项，就得到$\ell^2/(2r^2)$与$-3\ell^2/(8r^2)$。两个精确式的比值为
\[\begin{aligned}
\frac{E_a}{E_e}&=\frac{2(1+u)^{3/2}}{(1-u)^2}\\
&=2\left[1+\frac32u+O(u^2)\right]
\left[1+2u+O(u^2)\right]
=2\left[1+\frac72u+O(u^2)\right].
\end{aligned}\]
有限距离的比值从大于$2$的一侧逼近理想值。轴上靠近正电荷时分母趋零，偶极近似的相对误差不再小；$r>\ell/2$只保证精确式的区域正确，远场近似还要求$r\gg\ell$。
\end{solution}

\begin{exercise}[偶极在不均匀场中]\label{ex:es1-h6}
永久偶极矩$\vct p$固定，外场在偶极尺度内缓变。由两点受力求合力与力矩；若外场无旋，写成势能形式。 见图~\ref{fig:es1-h6-q}。
\end{exercise}

\begin{figure}[H]\centering
\begin{tikzpicture}[>=Stealth,line cap=round,line join=round]
\draw[axis](-2.7,0)--(2.9,0)node[right]{$x$};\fill(-0.8,0)circle(1.6pt)node[below]{$-q$};\fill(0.8,0)circle(1.6pt)node[below]{$+q$};\draw[axis](-.8,0)--(-1.4,-.6)node[left]{$-q\vct E_-$};\draw[axis](.8,0)--(1.7,.8)node[right]{$q\vct E_+$};
\end{tikzpicture}
\caption{负电荷受力与当地外场反向，合力来自两处场值之差。}\label{fig:es1-h6-q}
\end{figure}
\begin{solution}
以中心为$\vct r$，两电荷在$\vct r\pm\vct\ell/2$。展开场差，
\[F_i=q[E_i(\vct r+\vct\ell/2)-E_i(\vct r-\vct\ell/2)]
=p_j\partial_jE_i+O(q\ell^3\|\nabla^3\vct E\|).\]
图~\ref{fig:es1-h6-s}中差分沿$\vct p$作用于每个电场分量。力矩的首项来自两支力对中心的力臂相加，为$\vct\tau=\vct p\times\vct E$。

在静电无旋场中，$\partial_jE_i=\partial_iE_j$，固定$\vct p$时
\[\boxed{\vct F=\nabla(\vct p\cdot\vct E),\qquad U=-\vct p\cdot\vct E}.\]
永久偶极的平移与转动分别对应位置与方向的变化。若偶极矩由外场感生，其能量还包含形成偶极的代价，不能先令$p=\alpha E$再把永久偶极势能原样使用。
\begin{figure}[H]\centering
\begin{tikzpicture}[>=Stealth,line cap=round,line join=round]
\node at(0,0){$E_i(\vct r-\vct\ell/2)$};\node at(5,0){$E_i(\vct r+\vct\ell/2)$};\draw[axis](1.8,.15)--(3.2,.15)node[midway,above]{$\vct\ell\cdot\nabla$};\node at(2.5,-.9){$F_i=p_j\partial_jE_i$};
\end{tikzpicture}
\caption{有限场差在偶极近似下成为沿偶极轴的方向导数。}\label{fig:es1-h6-s}
\end{figure}


力矩可写为$(\vct\ell/2)\times q\vct E_++(-\vct\ell/2)\times(-q\vct E_-)$。保留均匀场部分给$q\vct\ell\times\vct E$，一阶梯度修正因两力臂对称而相消。若电场均匀，合力为零而力矩可不为零；若偶极已与外场对齐，力矩可为零而梯度力仍存在。力矩取决于局部电场与偶极方向，合力还取决于场的空间导数。
\end{solution}

\begin{exercise}[两偶极的角度与作用能]\label{ex:es1-h7}
两个永久偶极矩$\vct p_1,\vct p_2$相距$R$，尺寸均远小于$R$。写出相互作用能，比较同轴同向、并排同向两种构形的吸引或排斥。 见图~\ref{fig:es1-h7-q}。
\end{exercise}

\begin{figure}[H]\centering
\begin{tikzpicture}[>=Stealth,line cap=round,line join=round]
\fill(0,0)circle(1.4pt);\fill(3,0)circle(1.4pt);\draw[axis](0,0)--(3,0)node[midway,below]{$R\vct n$};\draw[axis](0,.25)--(.7,.95)node[above]{$\vct p_1$};\draw[axis](3,.25)--(3.3,1.2)node[above]{$\vct p_2$};
\end{tikzpicture}
\caption{连心单位矢量从第一个偶极指向第二个，偶极方向独立给定。}\label{fig:es1-h7-q}
\end{figure}
\begin{solution}
第一个偶极在第二个中心产生的场为$\vct E_1=\ke[3(\vct p_1\cdot\vct n)\vct n-\vct p_1]/R^3$。将第二个偶极的势能写成$-\vct p_2\cdot\vct E_1$，得
\[\boxed{U=\frac{\ke}{R^3}[\vct p_1\cdot\vct p_2
-3(\vct p_1\cdot\vct n)(\vct p_2\cdot\vct n)]}.\]
同轴同向时$U=-2\ke p_1p_2/R^3$，固定方向的径向力$-\dd U/\dd R=-6\ke p_1p_2/R^4$，为吸引；并排同向时$U=\ke p_1p_2/R^3$，径向力为$+3\ke p_1p_2/R^4$，为排斥。

图~\ref{fig:es1-h7-s}将两种几何分开。若方向没有约束，还须同时计算力矩，体系可能先发生转动，不能只由固定方向的径向力断言完整平衡稳定。偶极场按$R^{-3}$衰减，偶极间力再多一次空间求导，按$R^{-4}$衰减。
\begin{figure}[H]\centering
\begin{tikzpicture}[>=Stealth,line cap=round,line join=round]
\draw[axis](0,0)--(1,0)node[above]{$\vct p_1$};\draw[axis](2,0)--(3,0)node[above]{$\vct p_2$};\node at(1.5,-.6){同轴同向：吸引};\begin{scope}[shift={(5,0)}]\draw[axis](0,0)--(0,1);\draw[axis](2,0)--(2,1);\node at(1,-.6){并排同向：排斥};\end{scope}
\end{tikzpicture}
\caption{相同偶极方向在不同连心几何下，作用能的符号不同。}\label{fig:es1-h7-s}
\end{figure}


一般方向还存在非径向力，因为$\vct n$随横向位移变化。令$c_i=\vct p_i\cdot\vct n$，固定两偶极方向时，$\partial_jR=n_j$、$\partial_jn_i=(\delta_{ij}-n_in_j)/R$，因此
\[\nabla c_i=\frac{\vct p_i-c_i\vct n}{R},\qquad
\nabla R^{-3}=-\frac{3\vct n}{R^4}.
\]
将$U=\ke R^{-3}(\vct p_1\cdot\vct p_2-3c_1c_2)$作乘积求导，
\[\begin{split}
-\nabla U={}&\frac{3\ke}{R^4}(\vct p_1\cdot\vct p_2-3c_1c_2)\vct n\\
&+\frac{3\ke}{R^4}[c_2(\vct p_1-c_1\vct n)+c_1(\vct p_2-c_2\vct n)].
\end{split}\]
合并径向的$-3c_1c_2$与另外两个$-c_1c_2$，得到
\[\vct F=\frac{3\ke}{R^4}[c_1\vct p_2+c_2\vct p_1+(\vct p_1\cdot\vct p_2)\vct n-5c_1c_2\vct n].\]
前面两种对称构形使横向项抵消，才可只用$-\dd U/\dd R$求完整合力。任意夹角若只求径向导数，得到的只是力在连心方向的投影。
\end{solution}

\begin{exercise}[受诱导的小球]\label{ex:es1-h8}
小中性物体的各向同性极化率为$\alpha>0$，置于缓变静电场中，$\vct p=\alpha\vct E$。求平衡偶极的能量和力，并与永久偶极比较。 见图~\ref{fig:es1-h8-q}。
\end{exercise}

\begin{figure}[H]\centering
\begin{tikzpicture}[>=Stealth,line cap=round,line join=round]
\draw[axis](0,0)--(4.8,0)node[right]{$x$};\draw[curve](0,.8)..controls(2,.6)and(3,.3)..(4.5,.15);\draw[curve](0,-.8)..controls(2,-.6)and(3,-.3)..(4.5,-.15);\draw(2.5,0)circle(.25);\node[above]at(2.5,.5){$\vct p=\alpha\vct E$};
\end{tikzpicture}
\caption{物体的偶极矩随当地场改变，需把形成极化的能量计入。}\label{fig:es1-h8-q}
\end{figure}
\begin{solution}
在线性可逆模型中，维持偶极大小$p'$所需的共轭场为$p'/\alpha$，形成偶极的内能为
\[U_{\rm int}=\int_0^p\frac{p'}\alpha\dd p'
=\left[\frac{p'^2}{2\alpha}\right]_0^p=\frac{p^2}{2\alpha}.\]
与外场的相互作用为$-\vct p\cdot\vct E$。将内能与相互作用相加并配方，
\[U(\vct p)=\frac{p^2}{2\alpha}-\vct p\cdot\vct E
=\frac{|\vct p-\alpha\vct E|^2}{2\alpha}-\frac{\alpha E^2}{2}.
\]
这里$\alpha>0$，所以第一项在$\vct p=\alpha\vct E$时取最小值零。对位置求负梯度时再保持材料参数$\alpha$不变，得到$\vct p=\alpha\vct E$及
\[\boxed{U_{\rm eq}=-\frac12\alpha E^2,\qquad
\vct F=\frac\alpha2\nabla E^2}.\]
图~\ref{fig:es1-h8-s}的抛物线最低点随$E$移动，平衡能量是最低值，不能将$p=\alpha E$直接代入永久偶极的$-pE$后求导。

也可先在给定瞬时$p$时用$F_i=p_j\partial_iE_j$求机械力，再代入$p_j=\alpha E_j$，得到$\alpha E_j\partial_iE_j=\alpha\partial_i(E^2)/2$，与能量法一致。结果指向场强平方增大的方向，适用于物体尺寸小于场变化尺度且无不可逆极化的范围。
\begin{figure}[H]\centering
\begin{tikzpicture}[>=Stealth,line cap=round,line join=round]
\begin{scope}[x=1.7cm,y=1.5cm]\draw[axis](-0.2,0)--(2.48,0)node[right]{$p/(\alpha E)$};\draw[axis](0,-0.7)--(0,1)node[above]{$U/(\alpha E^2)$};\draw[curve,domain=-0.2:2.3,samples=101]plot(\x,{.5*(\x)^2-\x});\draw[guide](1,0)--(1,-.5);\fill(1,-.5)circle(1.5pt);\node[below]at(1,-.5){平衡};\end{scope}
\end{tikzpicture}
\caption{形成极化的正二次能与外场作用能相加，最低能量为$-\alpha E^2/2$。}\label{fig:es1-h8-s}
\end{figure}


若沿$x$方向的场大小为$E(x)$，则$F_x=\alpha E E'$。例如远处点电荷场$E=\ke Q/x^2$，有$F_x=-2\alpha\ke^2Q^2/x^5$，始终指向点电荷而与$Q$的正负无关。诱导物体虽然中性，近端与远端电力不等仍产生合力；这一$x^{-5}$行为也对应中性导体球像法中远处两个低阶力项相消后的首项。
\end{solution}

\begin{exercise}[柱坐标中的径向与绕轴场]\label{ex:es1-h9}
在$s>0$区域，$\vct A=a\vct e_s/s+b\vct e_\phi/s$。求散度、旋度，并计算单位长度圆柱侧面通量与绕轴一周环量。 见图~\ref{fig:es1-h9-q}。
\end{exercise}

\begin{figure}[H]\centering
\begin{tikzpicture}[>=Stealth,line cap=round,line join=round]
\draw[curve](0,0)circle(1.5);\draw[axis](0,0)--(1.5,0)node[midway,below]{$s$};\draw[axis](1.5,0)--(2.2,0)node[right]{$a/s$};\draw[axis](1.5,0)--(1.5,.8)node[above]{$b/s$};
\end{tikzpicture}
\caption{径向分量穿过圆柱侧壁，角向分量沿绕轴路径。}\label{fig:es1-h9-q}
\end{figure}
\begin{solution}
柱坐标中$\nabla\cdot\vct A=s^{-1}\partial_s(sA_s)+s^{-1}\partial_\phi A_\phi+\partial_zA_z=0$。轴外各旋度分量也为零，尤其$(\nabla\times\vct A)_z=s^{-1}\partial_s(sA_\phi)=0$。

图~\ref{fig:es1-h9-s}使用半径$s$、长度$L$的圆柱，侧通量为$2\pi sL(a/s)=2\pi aL$，端面通量为零。绕轴环量则为$2\pi s(b/s)=2\pi b$。局部散度和旋度为零，与这些有限积分不矛盾，因为闭面内部或回路所跨曲面包含了轴上的奇异源。

若改在不包围轴的可收缩小回路上积分，环量为零。判断能否使用全域积分定理，必须同时检查场在所围区域内是否光滑；不能只检查积分路径本身。
\begin{figure}[H]\centering
\begin{tikzpicture}[>=Stealth,line cap=round,line join=round]
\draw[curve](-1.4,0)--(-1.4,2) (1.4,0)--(1.4,2);\draw(0,0)ellipse(1.4 and .35);\draw(0,2)ellipse(1.4 and .35);\draw[guide](0,-.5)--(0,2.5);\draw[axis](1.4,1)--(2.2,1)node[right]{$\vct n$};\node[left]at(-1.4,1){$L$};
\end{tikzpicture}
\caption{轴线是被排除的奇异集，圆柱通量与绕轴环量分别测量两种源。}\label{fig:es1-h9-s}
\end{figure}


若用半径$s_1,s_2$围成不含轴的环柱，外侧与内侧有向通量分别为$2\pi aL$、$-2\pi aL$，总和为零。相同道理，环域边界的两条相反方向圆周环量相加为零。因此积分定理在光滑环域上完全成立，非零的单个圆周量并不构成反例。需要排除的是把穿过奇轴的实心面或体积当成光滑积分域。
\end{solution}

\begin{exercise}[绕轴场的势与分支]\label{ex:es1-h10}
在去掉$z$轴的空间，$\vct E=b\vct e_\phi/s$。它局部无旋，是否存在全域单值静电势？若切去一张半平面，写出局部势。 见图~\ref{fig:es1-h10-q}。
\end{exercise}

\begin{figure}[H]\centering
\begin{tikzpicture}[>=Stealth,line cap=round,line join=round]
\draw[curve](0,0)circle(1.5);\draw[axis](1.5,0)arc(0:50:1.5);\fill(0,0)circle(1.5pt);\draw[thick](0,0)--(2,0);\node[above]at(1,.05){切面};
\end{tikzpicture}
\caption{闭路能绕过被挖去的轴，局部无旋不保证绕轴环量为零。}\label{fig:es1-h10-q}
\end{figure}
\begin{solution}
绕轴一周的环量为$2\pi b$。当$b\ne0$，若存在单值势，闭路积分应为零，故矛盾；这个场不能作为该全域内的静电场。

切去图~\ref{fig:es1-h10-s}的半平面并规定$0<\phi<2\pi$后，可定义$\varphi=-b\phi+C$，因为$E_\phi=-(1/s)\partial_\phi\varphi=b/s$。两侧极限相差$2\pi b$，所以它只是所选切域上的单值势。切面位置可改变，绕轴的非零环量不变。

本例中的对数与角度常可合成复对数的实、虚部，但物理判断仍应回到环量和定义域。线电荷的径向$1/s$场有单值对数势，绕轴的角向$1/s$场则具有上述角度分支，两者不可混为一谈。
\begin{figure}[H]\centering
\begin{tikzpicture}[>=Stealth,line cap=round,line join=round]
\draw[curve](0,0)circle(1.5);\draw[thick](0,0)--(2,0);\node[above]at(1.2,.1){$\phi=0^+$};\node[below]at(1.2,-.1){$\phi=2\pi^-$};\node[below]at(0,-1.9){$\Delta\varphi=-2\pi b$};
\end{tikzpicture}
\caption{切面两侧的势极限不同，反映全域单值势不存在。}\label{fig:es1-h10-s}
\end{figure}


局部势的梯度在切面两侧都给同一个角向场，势值跳跃却不能当作真实静电偶极面的势跳跃：这里切面是人为选取的坐标分支，并没有题设给出的物理源。若全空间出现这种绕轴电场，通常须有随时间变化的穿轴磁通来维持非零环量，已经超出静电假设。局部数学表示可以保留，物理解释必须与题设的时间条件一致。
\end{solution}

\begin{exercise}[从场反求电荷]\label{ex:es1-h11}
球对称电场为$E_r=Ar$（$r<a$），$E_r=B/r^2$（$r>a$），$A>0$。求体电荷密度与球面电荷密度，并求无穷远零势。 见图~\ref{fig:es1-h11-q}。
\end{exercise}

\begin{figure}[H]\centering
\begin{tikzpicture}[>=Stealth,line cap=round,line join=round]
\draw[curve](0,0)circle(1.5);\fill(0,0)circle(1.4pt)node[below left]{$O$};\draw[axis](0,0)--(1.5,0)node[midway,above]{$a$};\node at(-.65,.7){$Ar$};\node[right]at(1.8,.8){$B/r^2$};
\end{tikzpicture}
\caption{内外两个场分支在球面可以连续，也可以有跃变。}\label{fig:es1-h11-q}
\end{figure}
\begin{solution}
球内$\partial_r(r^2Ar)=\partial_r(Ar^3)=3Ar^2$，所以$\rho=\epsilon_0r^{-2}(3Ar^2)=3\epsilon_0A$；球外$r^2E_r=B$的导数为零，体密度为零。球面法向场跃变给出
\[\boxed{\sigma=\epsilon_0(B/a^2-Aa)}.\]
图~\ref{fig:es1-h11-s}将两侧极限分开标记，只有$B=Aa^3$时面电荷消失。总电荷为$4\pi\epsilon_0B$，等于体电荷与面电荷之和。

由外向内积分，$\varphi_{\rm out}=-B\int_\infty^r s^{-2}\dd s=-B[-s^{-1}]_\infty^r=B/r$，球内
\[\varphi_{\rm in}=B/a+\int_r^aAs\dd s
=B/a+\left[\frac{As^2}{2}\right]_r^a
=B/a+A(a^2-r^2)/2.\]
势连续，导数的跃变恰好恢复面电荷。若只对两个开区域分别取散度，会漏掉球面的集中电荷。
\begin{figure}[H]\centering
\begin{tikzpicture}[>=Stealth,line cap=round,line join=round]
\draw[axis](0,0)--(3.5,0)node[right]{$r$};\draw[axis](0,0)--(0,2.4)node[above]{$E_r$};\draw[curve](0,0)--(1.5,1.2);\draw[curve,domain=1.5:3.2]plot(\x,{4.05/((\x)^2)});\draw[guide](1.5,0)--(1.5,1.8);\node[below]at(1.5,0){$a$};\node[left]at(1.5,1.2){$Aa$};\node[right]at(1.5,1.8){$B/a^2$};
\end{tikzpicture}
\caption{图示外侧场大于内侧场，因而球面电荷为正。}\label{fig:es1-h11-s}
\end{figure}


将面电荷项计入Gauss电量，$Q_V=4\pi\epsilon_0Aa^3$，$Q_S=4\pi\epsilon_0(B-Aa^3)$，相加恰为$4\pi\epsilon_0B$。当$B< Aa^3$，外层面电荷为负，甚至可使总电荷为零；例如$B=0$时球外场消失而内部仍有非零场。这样的屏蔽来自正体电荷与负表面电荷的共同作用，不能由球外零场推断内部电荷处处为零。
\end{solution}

\begin{exercise}[线偶极的远区场]\label{ex:es1-h12}
两条异号等量无限线电荷相距$2l$，线密度为$\pm\lambda$。在远区$r\gg l$求电势首项和场线，角度从由负线指向正线的轴量起。 见图~\ref{fig:es1-h12-q}。
\end{exercise}

\begin{figure}[H]\centering
\begin{tikzpicture}[>=Stealth,line cap=round,line join=round]
\draw[axis](-2.7,0)--(2.9,0)node[right]{$x$};\fill(-1.2,0)circle(1.6pt)node[below]{$-\lambda$};\fill(1.2,0)circle(1.6pt)node[below]{$+\lambda$};\fill(0.6,1.6)circle(1.5pt)node[above right]{$P$};\draw[guide](-1.2,0)--(0.6,1.6);\draw[guide](1.2,0)--(0.6,1.6);\draw[axis](-1.2,-.6)--(1.2,-.6)node[midway,below]{$d=2\lambda l$};
\end{tikzpicture}
\caption{这里的$d$是单位轴长的偶极矩，量纲与点偶极矩不同。}\label{fig:es1-h12-q}
\end{figure}
\begin{solution}
二维对数势的展开在本解答末尾给出。其首项为$\varphi=d\cos\theta/(2\pi\epsilon_0r)$，$d=2\lambda l$。分别求导，
\[E_r=-\frac{d\cos\theta}{2\pi\epsilon_0}(-r^{-2}),\qquad
E_\theta=-\frac1r\frac{d}{2\pi\epsilon_0r}(-\sin\theta),\]
即
\[E_r=\frac{d\cos\theta}{2\pi\epsilon_0r^2},\qquad
E_\theta=\frac{d\sin\theta}{2\pi\epsilon_0r^2}.\]
切线方程为$\dd r/r=\cot\theta\dd\theta$。令$u=\sin\theta$，$\dd u=\cos\theta\dd\theta$，则$\ln r=\int\dd u/u+C_0=\ln|\sin\theta|+C_0$，故$\boxed{r=C\sin\theta}$，常数的符号随半平面选择。图~\ref{fig:es1-h12-s}中该式是通过原点、与偶极轴相切的圆族在所选坐标下的表示。

点偶极的势按$r^{-2}$衰减，线偶极的势按$r^{-1}$衰减；相应场线分别含$\sin^2\theta$与$\sin\theta$。幂次差来自二维对数核与三维库仑核的差别，不能把两种偶极只改一个系数便互相使用。
\begin{figure}[H]\centering
\begin{tikzpicture}[>=Stealth,line cap=round,line join=round]
\draw[axis](-2.2,0)--(2.4,0)node[right]{$x$};\draw[axis](0,-1.6)--(0,1.8)node[above]{$y$};\draw[curve](0,.8)circle(.8);\draw[curve](0,-.8)circle(.8);\fill(0,0)circle(1.5pt);\node[right]at(2.6,.5){极轴取$x$方向};
\end{tikzpicture}
\caption{以$x$为极轴，$r=C\sin\theta$是一族圆；$C$的正负选择对应上下半平面。}\label{fig:es1-h12-s}
\end{figure}


从两个对数势直接展开，正、负线源分别在$x=+l,-l$，对应距离满足$r_\pm^2=r^2(1\mp2t\cos\theta+t^2)$，其中$t=l/r$。记$c=\cos\theta$，使用$\ln(1+v)=v-v^2/2+O(v^3)$，
\[\begin{aligned}
\ln\frac{r_\pm}r
&=\frac12\ln(1\mp2tc+t^2)\\
&=\frac12\left[\mp2tc+t^2-2t^2c^2+O(t^3)\right]\\
&=\mp tc+\left(\frac12-c^2\right)t^2+O(t^3).
\end{aligned}\]
两式相减，二次项抵消，得到
\[\ln\frac{r_-}{r_+}=2t\cos\theta+O(t^3),\qquad
\varphi=\frac{2\lambda l}{2\pi\epsilon_0r}\cos\theta
+O\left(\frac{\lambda l^3}{\epsilon_0r^3}\right).
\]
所以$d=2\lambda l$，方向仍由负线源指向正线源。极坐标曲线$r=C\sin\theta$乘以$r$后，逐步整理为
\[r^2=Cr\sin\theta,\quad x^2+y^2=Cy,
\quad x^2+\left(y-\frac C2\right)^2=\frac{C^2}{4}.
\]
圆心为$(0,C/2)$、半径为$|C|/2$，上下两圆都是线偶极极限的精确场线。原点属于理想源区，不能因圆通过该点而把它当作普通的光滑场点。\end{solution}

\begin{exercise}[四极矩的尺度]\label{ex:es1-h13}
保持四极参数$A=q\ell^2$不变，令$\ell\to0$。比较电荷$q$、偶极势与四极势的尺度，并判断该理想化能否用于源点处的能量。 见图~\ref{fig:es1-h13-q}。
\end{exercise}

\begin{figure}[H]\centering
\begin{tikzpicture}[>=Stealth,line cap=round,line join=round]
\draw[axis](-2.7,0)--(2.9,0)node[right]{$x$};\fill(-1.2,0)circle(1.6pt)node[below]{$-q$};\fill(0,0)circle(1.6pt)node[below]{$+2q$};\fill(1.2,0)circle(1.6pt)node[below]{$-q$};\fill(0.4,1.8)circle(1.5pt)node[above right]{$P$};\draw[guide](-1.2,0)--(0.4,1.8);\draw[guide](0,0)--(0.4,1.8);\draw[guide](1.2,0)--(0.4,1.8);
\end{tikzpicture}
\caption{四极极限同时缩小几何尺度并增大电荷，观察点保持在源区之外。}\label{fig:es1-h13-q}
\end{figure}
\begin{solution}
固定$A$意味着$q=A/\ell^2$，不能把$q$也当作固定常数。总电荷与偶极矩仍严格为零，四极势首项$-\ke A(3\cos^2\theta-1)/r^3$在固定$r>0$处保持有限，误差为$O(\ke A\ell^2/r^5)$并趋零。

图~\ref{fig:es1-h13-s}的观察点与源尺度之比趋于无穷，保证展开合法。但源电荷间相互作用能量按$\ke q^2/\ell\sim\ke A^2/\ell^5$发散，理想极限不能给出有限的自身能。外场中的相互作用可用适当多极矩描述，自能则需保留源的有限结构。

同样，理想偶极极限固定$p=q\ell$时$q\sim\ell^{-1}$。偶极极限固定$p$，四极极限固定$A$；只规定距离趋零，尚未确定相应的源模型。
\begin{figure}[H]\centering
\begin{tikzpicture}[>=Stealth,line cap=round,line join=round]
\draw[curve](-.5,-.5)rectangle(.5,.5);\node[below]at(0,-.65){源区$O(\ell)$};\fill(3,1.3)circle(1.5pt)node[right]{$P$};\draw[axis](0,0)--(3,1.3)node[midway,above]{$r$固定};\node[right]at(2,-.65){$\ell/r\to0$};
\end{tikzpicture}
\caption{远区极限控制的是固定观察点的外场，并未控制源内部的相互作用能。}\label{fig:es1-h13-s}
\end{figure}


三个有限源的互能也可算出发散系数。两对相邻异号电荷各给$-2\ke q^2/\ell$，两端同号负电荷给$+\ke q^2/(2\ell)$，故不计各点自能的总互能为$-7\ke q^2/(2\ell)=-7\ke A^2/(2\ell^5)$。它向负无穷发散，仍不能与点电荷本来就无穷的自能随意相加抵消；连续有限源模型才可讨论完整可测的自身能。
\end{solution}

\begin{exercise}[势的角向因子与场的衰减]\label{ex:es1-h14}
轴对称试探势$\varphi=C r^{-n}\cos\theta$，$n>0$。求它在$r>0$的电荷密度；何时可以作为无源偶极外势？ 见图~\ref{fig:es1-h14-q}。
\end{exercise}

\begin{figure}[H]\centering
\begin{tikzpicture}[>=Stealth,line cap=round,line join=round]
\draw[curve](0,0)circle(1.5);\fill(0,0)circle(1.4pt)node[below left]{$O$};\draw[axis](0,0)--(1.2,.9)node[right]{$r$};\draw[axis](0,0)--(0,2)node[above]{$z$};\node[left]at(-1.7,.8){$Cr^{-n}\cos\theta$};
\end{tikzpicture}
\caption{径向幂次与角函数不能独立指定，还需满足无源场方程。}\label{fig:es1-h14-q}
\end{figure}
\begin{solution}
径向部分先有$\partial_r(r^{-n})=-nr^{-n-1}$，再乘$r^2$得$-nr^{1-n}$，故
\[r^{-2}\partial_r[-nr^{1-n}]=r^{-2}[-n(1-n)r^{-n}]=n(n-1)r^{-n-2}.\]
角向部分可逐步计算为
\[\frac1{\sin\theta}\frac{\partial}{\partial\theta}
\left(\sin\theta\frac{\partial\cos\theta}{\partial\theta}\right)
=\frac1{\sin\theta}\frac{\partial(-\sin^2\theta)}{\partial\theta}
=-2\cos\theta.
\]
它在Laplace算子中还要乘$1/r^2$，故
\[\nabla^2\varphi=C[n(n-1)-2]r^{-n-2}\cos\theta,
\quad\rho=-\epsilon_0\nabla^2\varphi.\]
要使源外密度为零，必须$n(n-1)=2$，解为$n=2$或$n=-1$。题设$n>0$只留下$n=2$，即偶极衰减；$n=-1$若放回一般问题，则是随$r$增长的匀强场势。

图~\ref{fig:es1-h14-s}用同一个角函数比较两种允许的径向行为。单靠量纲或“远处递减”不能任意选择幂次，Laplace方程同时约束径向和角向部分。
\begin{figure}[H]\centering
\begin{tikzpicture}[>=Stealth,line cap=round,line join=round]
\draw[axis](0,0)--(4,0)node[right]{$r$};\draw[axis](0,0)--(0,2.4)node[above]{径向因子};\draw[curve,domain=.8:3.5]plot(\x,{1/((\x)^2)});\draw[dashed](0,0)--(3.5,1.75);\node[right]at(4,1.7){虚线：$r$};\node[right]at(4,.5){实线：$r^{-2}$};
\end{tikzpicture}
\caption{同一一阶角函数的两个无源径向解是$r$与$r^{-2}$，边界条件决定保留哪一个。}\label{fig:es1-h14-s}
\end{figure}


若允许任意球谐$Y_{lm}$，角算子的本征值为$-l(l+1)$，径向幂$r^s$满足$s(s+1)=l(l+1)$，根为$s=l,-l-1$。本题是$l=1$，因此正则内解$r$与衰减外解$r^{-2}$成一对。这个规律是第二讲分离变量的三维版本，也解释了为何四极势必须按$r^{-3}$而不是任意负幂衰减。
\end{solution}

\begin{exercise}[一条场线上的功]\label{ex:es1-h15}
理想偶极场中，一条场线$r=C\sin^2\theta$上取$P_1,P_2$，$0<\theta_1<\theta_2<\pi/2$。求正试探电荷$q_0$沿该场线移动时电场所做的功，并判定符号。 见图~\ref{fig:es1-h15-q}。
\end{exercise}

\begin{figure}[H]\centering
\begin{tikzpicture}[>=Stealth,line cap=round,line join=round]
\draw[axis](-2.3,0)--(2.45,0)node[right]{$x$};\draw[axis](0,-2.3)--(0,2.45)node[above]{$z$};\draw[curve]plot coordinates{(-0.000000,0.000001) (-0.000034,0.001049) (-0.000259,0.004056) (-0.000859,0.008992) (-0.002012,0.015811) (-0.003894,0.024442) (-0.006671,0.034803) (-0.010501,0.046789) (-0.015532,0.060283) (-0.021900,0.075150) (-0.029727,0.091242) (-0.039121,0.108398) (-0.050174,0.126445) (-0.062963,0.145203) (-0.077543,0.164481) (-0.093957,0.184082) (-0.112224,0.203806) (-0.132346,0.223449) (-0.154307,0.242806) (-0.178070,0.261672) (-0.203578,0.279848) (-0.230758,0.297136) (-0.259518,0.313345) (-0.289745,0.328294) (-0.321315,0.341809) (-0.354084,0.353730) (-0.387894,0.363907) (-0.422576,0.372206) (-0.457946,0.378507) (-0.493811,0.382707) (-0.529970,0.384722) (-0.566213,0.384483) (-0.602325,0.381943) (-0.638088,0.377072) (-0.673283,0.369860) (-0.707688,0.360318) (-0.741086,0.348475) (-0.773261,0.334381) (-0.804006,0.318105) (-0.833119,0.299734) (-0.860406,0.279373) (-0.885687,0.257143) (-0.908791,0.233183) (-0.929564,0.207645) (-0.947862,0.180697) (-0.963564,0.152515) (-0.976560,0.123289) (-0.986762,0.093217) (-0.994099,0.062504) (-0.998522,0.031360) (-1.000000,0.000000) (-0.998522,-0.031360) (-0.994099,-0.062504) (-0.986762,-0.093217) (-0.976560,-0.123289) (-0.963564,-0.152515) (-0.947862,-0.180697) (-0.929564,-0.207645) (-0.908791,-0.233183) (-0.885687,-0.257143) (-0.860406,-0.279373) (-0.833119,-0.299734) (-0.804006,-0.318105) (-0.773261,-0.334381) (-0.741086,-0.348475) (-0.707688,-0.360318) (-0.673283,-0.369860) (-0.638088,-0.377072) (-0.602325,-0.381943) (-0.566213,-0.384483) (-0.529970,-0.384722) (-0.493811,-0.382707) (-0.457946,-0.378507) (-0.422576,-0.372206) (-0.387894,-0.363907) (-0.354084,-0.353730) (-0.321315,-0.341809) (-0.289745,-0.328294) (-0.259518,-0.313345) (-0.230758,-0.297136) (-0.203578,-0.279848) (-0.178070,-0.261672) (-0.154307,-0.242806) (-0.132346,-0.223449) (-0.112224,-0.203806) (-0.093957,-0.184082) (-0.077543,-0.164481) (-0.062963,-0.145203) (-0.050174,-0.126445) (-0.039121,-0.108398) (-0.029727,-0.091242) (-0.021900,-0.075150) (-0.015532,-0.060283) (-0.010501,-0.046789) (-0.006671,-0.034803) (-0.003894,-0.024442) (-0.002012,-0.015811) (-0.000859,-0.008992) (-0.000259,-0.004056) (-0.000034,-0.001049) (-0.000000,-0.000001)};\draw[curve]plot coordinates{(0.000000,0.000001) (0.000034,0.001049) (0.000259,0.004056) (0.000859,0.008992) (0.002012,0.015811) (0.003894,0.024442) (0.006671,0.034803) (0.010501,0.046789) (0.015532,0.060283) (0.021900,0.075150) (0.029727,0.091242) (0.039121,0.108398) (0.050174,0.126445) (0.062963,0.145203) (0.077543,0.164481) (0.093957,0.184082) (0.112224,0.203806) (0.132346,0.223449) (0.154307,0.242806) (0.178070,0.261672) (0.203578,0.279848) (0.230758,0.297136) (0.259518,0.313345) (0.289745,0.328294) (0.321315,0.341809) (0.354084,0.353730) (0.387894,0.363907) (0.422576,0.372206) (0.457946,0.378507) (0.493811,0.382707) (0.529970,0.384722) (0.566213,0.384483) (0.602325,0.381943) (0.638088,0.377072) (0.673283,0.369860) (0.707688,0.360318) (0.741086,0.348475) (0.773261,0.334381) (0.804006,0.318105) (0.833119,0.299734) (0.860406,0.279373) (0.885687,0.257143) (0.908791,0.233183) (0.929564,0.207645) (0.947862,0.180697) (0.963564,0.152515) (0.976560,0.123289) (0.986762,0.093217) (0.994099,0.062504) (0.998522,0.031360) (1.000000,0.000000) (0.998522,-0.031360) (0.994099,-0.062504) (0.986762,-0.093217) (0.976560,-0.123289) (0.963564,-0.152515) (0.947862,-0.180697) (0.929564,-0.207645) (0.908791,-0.233183) (0.885687,-0.257143) (0.860406,-0.279373) (0.833119,-0.299734) (0.804006,-0.318105) (0.773261,-0.334381) (0.741086,-0.348475) (0.707688,-0.360318) (0.673283,-0.369860) (0.638088,-0.377072) (0.602325,-0.381943) (0.566213,-0.384483) (0.529970,-0.384722) (0.493811,-0.382707) (0.457946,-0.378507) (0.422576,-0.372206) (0.387894,-0.363907) (0.354084,-0.353730) (0.321315,-0.341809) (0.289745,-0.328294) (0.259518,-0.313345) (0.230758,-0.297136) (0.203578,-0.279848) (0.178070,-0.261672) (0.154307,-0.242806) (0.132346,-0.223449) (0.112224,-0.203806) (0.093957,-0.184082) (0.077543,-0.164481) (0.062963,-0.145203) (0.050174,-0.126445) (0.039121,-0.108398) (0.029727,-0.091242) (0.021900,-0.075150) (0.015532,-0.060283) (0.010501,-0.046789) (0.006671,-0.034803) (0.003894,-0.024442) (0.002012,-0.015811) (0.000859,-0.008992) (0.000259,-0.004056) (0.000034,-0.001049) (0.000000,-0.000001)};\draw[curve]plot coordinates{(-0.000000,0.000002) (-0.000054,0.001678) (-0.000415,0.006489) (-0.001374,0.014388) (-0.003219,0.025297) (-0.006230,0.039108) (-0.010673,0.055685) (-0.016801,0.074863) (-0.024851,0.096453) (-0.035040,0.120240) (-0.047563,0.145987) (-0.062594,0.173436) (-0.080279,0.202313) (-0.100740,0.232325) (-0.124070,0.263169) (-0.150331,0.294531) (-0.179558,0.326090) (-0.211754,0.357518) (-0.246892,0.388489) (-0.284912,0.418676) (-0.325725,0.447757) (-0.369213,0.475417) (-0.415228,0.501353) (-0.463593,0.525271) (-0.514104,0.546895) (-0.566534,0.565968) (-0.620631,0.582251) (-0.676122,0.595529) (-0.732714,0.605611) (-0.790098,0.612332) (-0.847952,0.615555) (-0.905941,0.615174) (-0.963720,0.611109) (-1.020942,0.603315) (-1.077252,0.591776) (-1.132301,0.576508) (-1.185737,0.557560) (-1.237218,0.535010) (-1.286410,0.508969) (-1.332990,0.479575) (-1.376650,0.446996) (-1.417099,0.411429) (-1.454066,0.373093) (-1.487302,0.332233) (-1.516580,0.289114) (-1.541702,0.244024) (-1.562496,0.197262) (-1.578819,0.149147) (-1.590559,0.100006) (-1.597636,0.050176) (-1.600000,0.000000) (-1.597636,-0.050176) (-1.590559,-0.100006) (-1.578819,-0.149147) (-1.562496,-0.197262) (-1.541702,-0.244024) (-1.516580,-0.289114) (-1.487302,-0.332233) (-1.454066,-0.373093) (-1.417099,-0.411429) (-1.376650,-0.446996) (-1.332990,-0.479575) (-1.286410,-0.508969) (-1.237218,-0.535010) (-1.185737,-0.557560) (-1.132301,-0.576508) (-1.077252,-0.591776) (-1.020942,-0.603315) (-0.963720,-0.611109) (-0.905941,-0.615174) (-0.847952,-0.615555) (-0.790098,-0.612332) (-0.732714,-0.605611) (-0.676122,-0.595529) (-0.620631,-0.582251) (-0.566534,-0.565968) (-0.514104,-0.546895) (-0.463593,-0.525271) (-0.415228,-0.501353) (-0.369213,-0.475417) (-0.325725,-0.447757) (-0.284912,-0.418676) (-0.246892,-0.388489) (-0.211754,-0.357518) (-0.179558,-0.326090) (-0.150331,-0.294531) (-0.124070,-0.263169) (-0.100740,-0.232325) (-0.080279,-0.202313) (-0.062594,-0.173436) (-0.047563,-0.145987) (-0.035040,-0.120240) (-0.024851,-0.096453) (-0.016801,-0.074863) (-0.010673,-0.055685) (-0.006230,-0.039108) (-0.003219,-0.025297) (-0.001374,-0.014388) (-0.000415,-0.006489) (-0.000054,-0.001678) (-0.000000,-0.000002)};\draw[curve]plot coordinates{(0.000000,0.000002) (0.000054,0.001678) (0.000415,0.006489) (0.001374,0.014388) (0.003219,0.025297) (0.006230,0.039108) (0.010673,0.055685) (0.016801,0.074863) (0.024851,0.096453) (0.035040,0.120240) (0.047563,0.145987) (0.062594,0.173436) (0.080279,0.202313) (0.100740,0.232325) (0.124070,0.263169) (0.150331,0.294531) (0.179558,0.326090) (0.211754,0.357518) (0.246892,0.388489) (0.284912,0.418676) (0.325725,0.447757) (0.369213,0.475417) (0.415228,0.501353) (0.463593,0.525271) (0.514104,0.546895) (0.566534,0.565968) (0.620631,0.582251) (0.676122,0.595529) (0.732714,0.605611) (0.790098,0.612332) (0.847952,0.615555) (0.905941,0.615174) (0.963720,0.611109) (1.020942,0.603315) (1.077252,0.591776) (1.132301,0.576508) (1.185737,0.557560) (1.237218,0.535010) (1.286410,0.508969) (1.332990,0.479575) (1.376650,0.446996) (1.417099,0.411429) (1.454066,0.373093) (1.487302,0.332233) (1.516580,0.289114) (1.541702,0.244024) (1.562496,0.197262) (1.578819,0.149147) (1.590559,0.100006) (1.597636,0.050176) (1.600000,0.000000) (1.597636,-0.050176) (1.590559,-0.100006) (1.578819,-0.149147) (1.562496,-0.197262) (1.541702,-0.244024) (1.516580,-0.289114) (1.487302,-0.332233) (1.454066,-0.373093) (1.417099,-0.411429) (1.376650,-0.446996) (1.332990,-0.479575) (1.286410,-0.508969) (1.237218,-0.535010) (1.185737,-0.557560) (1.132301,-0.576508) (1.077252,-0.591776) (1.020942,-0.603315) (0.963720,-0.611109) (0.905941,-0.615174) (0.847952,-0.615555) (0.790098,-0.612332) (0.732714,-0.605611) (0.676122,-0.595529) (0.620631,-0.582251) (0.566534,-0.565968) (0.514104,-0.546895) (0.463593,-0.525271) (0.415228,-0.501353) (0.369213,-0.475417) (0.325725,-0.447757) (0.284912,-0.418676) (0.246892,-0.388489) (0.211754,-0.357518) (0.179558,-0.326090) (0.150331,-0.294531) (0.124070,-0.263169) (0.100740,-0.232325) (0.080279,-0.202313) (0.062594,-0.173436) (0.047563,-0.145987) (0.035040,-0.120240) (0.024851,-0.096453) (0.016801,-0.074863) (0.010673,-0.055685) (0.006230,-0.039108) (0.003219,-0.025297) (0.001374,-0.014388) (0.000415,-0.006489) (0.000054,-0.001678) (0.000000,-0.000002)};\draw[curve]plot coordinates{(-0.000000,0.000002) (-0.000075,0.002307) (-0.000570,0.008922) (-0.001889,0.019783) (-0.004427,0.034783) (-0.008566,0.053773) (-0.014676,0.076566) (-0.023102,0.102937) (-0.034171,0.132623) (-0.048180,0.165330) (-0.065399,0.200732) (-0.086066,0.238475) (-0.110384,0.278180) (-0.138518,0.319447) (-0.170596,0.361858) (-0.206705,0.404981) (-0.246893,0.448373) (-0.291162,0.491588) (-0.339476,0.534172) (-0.391753,0.575679) (-0.447872,0.615666) (-0.507668,0.653699) (-0.570939,0.689360) (-0.637440,0.722247) (-0.706893,0.751981) (-0.778984,0.778206) (-0.853368,0.800595) (-0.929667,0.818852) (-1.007481,0.832715) (-1.086385,0.841956) (-1.165934,0.846389) (-1.245668,0.845864) (-1.325115,0.840275) (-1.403795,0.829558) (-1.481222,0.813692) (-1.556913,0.792699) (-1.630388,0.766645) (-1.701175,0.735639) (-1.768814,0.699832) (-1.832861,0.659415) (-1.892894,0.614620) (-1.948512,0.565715) (-1.999341,0.513003) (-2.045040,0.456820) (-2.085297,0.397532) (-2.119840,0.335532) (-2.148431,0.271235) (-2.170876,0.205077) (-2.187019,0.137508) (-2.196749,0.068992) (-2.200000,0.000000) (-2.196749,-0.068992) (-2.187019,-0.137508) (-2.170876,-0.205077) (-2.148431,-0.271235) (-2.119840,-0.335532) (-2.085297,-0.397532) (-2.045040,-0.456820) (-1.999341,-0.513003) (-1.948512,-0.565715) (-1.892894,-0.614620) (-1.832861,-0.659415) (-1.768814,-0.699832) (-1.701175,-0.735639) (-1.630388,-0.766645) (-1.556913,-0.792699) (-1.481222,-0.813692) (-1.403795,-0.829558) (-1.325115,-0.840275) (-1.245668,-0.845864) (-1.165934,-0.846389) (-1.086385,-0.841956) (-1.007481,-0.832715) (-0.929667,-0.818852) (-0.853368,-0.800595) (-0.778984,-0.778206) (-0.706893,-0.751981) (-0.637440,-0.722247) (-0.570939,-0.689360) (-0.507668,-0.653699) (-0.447872,-0.615666) (-0.391753,-0.575679) (-0.339476,-0.534172) (-0.291162,-0.491588) (-0.246893,-0.448373) (-0.206705,-0.404981) (-0.170596,-0.361858) (-0.138518,-0.319447) (-0.110384,-0.278180) (-0.086066,-0.238475) (-0.065399,-0.200732) (-0.048180,-0.165330) (-0.034171,-0.132623) (-0.023102,-0.102937) (-0.014676,-0.076566) (-0.008566,-0.053773) (-0.004427,-0.034783) (-0.001889,-0.019783) (-0.000570,-0.008922) (-0.000075,-0.002307) (-0.000000,-0.000002)};\draw[curve]plot coordinates{(0.000000,0.000002) (0.000075,0.002307) (0.000570,0.008922) (0.001889,0.019783) (0.004427,0.034783) (0.008566,0.053773) (0.014676,0.076566) (0.023102,0.102937) (0.034171,0.132623) (0.048180,0.165330) (0.065399,0.200732) (0.086066,0.238475) (0.110384,0.278180) (0.138518,0.319447) (0.170596,0.361858) (0.206705,0.404981) (0.246893,0.448373) (0.291162,0.491588) (0.339476,0.534172) (0.391753,0.575679) (0.447872,0.615666) (0.507668,0.653699) (0.570939,0.689360) (0.637440,0.722247) (0.706893,0.751981) (0.778984,0.778206) (0.853368,0.800595) (0.929667,0.818852) (1.007481,0.832715) (1.086385,0.841956) (1.165934,0.846389) (1.245668,0.845864) (1.325115,0.840275) (1.403795,0.829558) (1.481222,0.813692) (1.556913,0.792699) (1.630388,0.766645) (1.701175,0.735639) (1.768814,0.699832) (1.832861,0.659415) (1.892894,0.614620) (1.948512,0.565715) (1.999341,0.513003) (2.045040,0.456820) (2.085297,0.397532) (2.119840,0.335532) (2.148431,0.271235) (2.170876,0.205077) (2.187019,0.137508) (2.196749,0.068992) (2.200000,0.000000) (2.196749,-0.068992) (2.187019,-0.137508) (2.170876,-0.205077) (2.148431,-0.271235) (2.119840,-0.335532) (2.085297,-0.397532) (2.045040,-0.456820) (1.999341,-0.513003) (1.948512,-0.565715) (1.892894,-0.614620) (1.832861,-0.659415) (1.768814,-0.699832) (1.701175,-0.735639) (1.630388,-0.766645) (1.556913,-0.792699) (1.481222,-0.813692) (1.403795,-0.829558) (1.325115,-0.840275) (1.245668,-0.845864) (1.165934,-0.846389) (1.086385,-0.841956) (1.007481,-0.832715) (0.929667,-0.818852) (0.853368,-0.800595) (0.778984,-0.778206) (0.706893,-0.751981) (0.637440,-0.722247) (0.570939,-0.689360) (0.507668,-0.653699) (0.447872,-0.615666) (0.391753,-0.575679) (0.339476,-0.534172) (0.291162,-0.491588) (0.246893,-0.448373) (0.206705,-0.404981) (0.170596,-0.361858) (0.138518,-0.319447) (0.110384,-0.278180) (0.086066,-0.238475) (0.065399,-0.200732) (0.048180,-0.165330) (0.034171,-0.132623) (0.023102,-0.102937) (0.014676,-0.076566) (0.008566,-0.053773) (0.004427,-0.034783) (0.001889,-0.019783) (0.000570,-0.008922) (0.000075,-0.002307) (0.000000,-0.000002)};\fill(0,0)circle(1.5pt);
\fill(0.584286,0.696324)circle(1.8pt)node[above]{$P_1$};\fill(1.637758,0.763699)circle(1.8pt)node[above]{$P_2$};
\end{tikzpicture}
\caption{两点在同一场线上，但电势通常不同。}\label{fig:es1-h15-q}
\end{figure}
\begin{solution}
电场功只与端点有关，
\[W=q_0[\varphi(P_1)-\varphi(P_2)]
=\frac{\ke pq_0}{C^2}
\left(\frac{\cos\theta_1}{\sin^4\theta_1}
-\frac{\cos\theta_2}{\sin^4\theta_2}\right).\]
图~\ref{fig:es1-h15-s}以 $\theta$ 为参数，北半球上 $E_\theta>0$，沿场方向的极角增大。对端点势中的角度因子求导，
\[\begin{aligned}
\frac{\dd}{\dd\theta}(\cos\theta\sin^{-4}\theta)
&=-\sin\theta\sin^{-4}\theta
-4\cos^2\theta\sin^{-5}\theta\\
&=-\sin^{-3}\theta-4\frac{\cos^2\theta}{\sin^5\theta}<0.
\end{aligned}\]
因此 $\varphi(P_1)>\varphi(P_2)$，$W>0$。

若沿反方向强制移动，电场做负功，外力在准静态过程中做等量正功。场线不是等势线，沿场线方向移动反而使势持续下降；正交关系只说明沿等势面的功为零。
\begin{figure}[H]\centering
\begin{tikzpicture}[>=Stealth,line cap=round,line join=round]
\begin{scope}[x=2cm,y=0.25cm]\draw[axis](0,0)--(1.68,0)node[right]{$\theta$};\draw[axis](0,0)--(0,9)node[above]{$C^2\varphi/(\ke p)$};\draw[curve,domain=0.6:1.5,samples=101]plot(\x,{cos(\x r)/((sin(\x r))^4)});\draw[axis](.8,3.5)--(1.1,1.4);\end{scope}

\end{tikzpicture}
\caption{北半球沿场线向增大的极角移动时，电势单调下降。}\label{fig:es1-h15-s}
\end{figure}


沿曲线直接积分也能恢复端点式。由$r=C\sin^2\theta$得$\dd r=2C\sin\theta\cos\theta\dd\theta$，电场微功为$q_0(E_r\dd r+E_\theta r\dd\theta)$。代入偶极分量后，积分核为$\ke pq_0[4\cos^2\theta/\sin^5\theta+1/\sin^3\theta]\dd\theta/C^2$，恰为$-\dd[\ke pq_0\cos\theta/(C^2\sin^4\theta)]$。两种算法一致，路径积分多出的几何因子来自球坐标线元。
\end{solution}

\begin{exercise}[点电荷加偶极后的场线]\label{ex:es1-h16}
总电荷$Q>0$和沿正$z$轴的理想偶极$p>0$位于同一原点。求轴对称场线的隐式式，判断是否有有限零场点，并说明远区的主导形状。 见图~\ref{fig:es1-h16-q}。
\end{exercise}

\begin{figure}[H]\centering
\begin{tikzpicture}[>=Stealth,line cap=round,line join=round]
\draw[axis](-1.8,0)--(2,0)node[right]{$x$};\draw[axis](0,-2)--(0,2.2)node[above]{$z$};\fill(0,0)circle(2pt)node[below right]{$Q$};\draw[axis](-.4,-.4)--(-.4,.6)node[left]{$\vct p$};
\end{tikzpicture}
\caption{单极与偶极位于同一展开原点，势与场逐项叠加。}\label{fig:es1-h16-q}
\end{figure}
\begin{solution}
势为$\varphi=\ke(Q/r+p\cos\theta/r^2)$。球冠通量直接相加，
\[\Psi=\frac Q{2\epsilon_0}(1-\cos\theta)
+\frac p{2\epsilon_0r}\sin^2\theta.
\]
所以场线满足$Q(1-\cos\theta)+(p/r)\sin^2\theta=C$。图~\ref{fig:es1-h16-s}中远处球面上，偶极通量项按$p/r$变小，场线逐渐接近单极的径向射线。

角向场$E_\theta=\ke p\sin\theta/r^3$只有在轴线上为零。正轴$E_r=\ke(Q/r^2+2p/r^3)>0$，无零点；负轴$E_r=\ke(Q/r^2-2p/r^3)$，在$r=2p/Q$为零。该点的直角坐标为$z=-2p/Q$，并非原点的有效替代源。

对于有限源体系，这个零点只有在$2p/Q$远大于源尺寸时才可由截断多极式可靠预测。若它落入源区，应回到精确分布重新求解。
\begin{figure}[H]\centering
\begin{tikzpicture}[>=Stealth,line cap=round,line join=round]
\draw[axis](-2,0)--(2.2,0)node[right]{$x$};\draw[axis](0,-2.4)--(0,2.4)node[above]{$z$};\fill(0,0)circle(1.5pt)node[right]{$Q,p$};\fill(0,-1.5)circle(1.6pt)node[right]{$z=-2p/Q$};\draw[guide](0,0)circle(2);\foreach\a in{20,60,120,160}{\draw[axis]({1.65*cos(\a)},{1.65*sin(\a)})--({2.1*cos(\a)},{2.1*sin(\a)});}
\end{tikzpicture}
\caption{负轴上的有限零场点与远处近径向场分开标示，零点处需另行分析场线分界。}\label{fig:es1-h16-s}
\end{figure}


在负轴$z<0$上，径向单位矢量为$-\vct e_z$，故轴向场须写为
\[E_z=-\frac{\ke Q}{z^2}-\frac{2\ke p}{z^3}.
\]
对$z$求导，再代零点$z_0=-2p/Q$，
\[\begin{aligned}
\partial_zE_z&=\frac{2\ke Q}{z^3}+\frac{6\ke p}{z^4},\\
\partial_zE_z(z_0)&=-\frac{\ke Q^4}{4p^3}
+\frac{3\ke Q^4}{8p^3}
=\frac{\ke Q^4}{8p^3}>0.
\end{aligned}\]
所以正试探电荷向上小移后，轴向电力也向上，沿轴失稳。零点附近无源且绕轴对称，两个横向场斜率相等，再由散度为零得
\[\partial_xE_x=\partial_yE_y
=-\frac12\partial_zE_z=-\frac{\ke Q^4}{16p^3}<0.
\]
因此横向有回复。约束在通过零点、垂直于轴的平面内时可以稳定，自由运动时则不稳定。该结论仍要求零点位于多极近似有效的源外区域。\end{solution}
